\chapter{Integration and Stokes formulae}
\index{integration!tropical}\index{Stokes formula!polyhedral}
\label{chap:tropical-integration}

Integration is where weights, lattice normalizations, and balancing become necessary.
We first prove Stokes on one polyhedron, then show that balancing cancels
internal faces. This is the mechanism behind the global CLD theorem.

Two chapters of preparation have produced a bigraded algebra of superforms
(\Cref{chap:superforms}) and a supply of supports on which such forms live
(\Cref{chap:tropical-spaces}).  Neither chapter has yet produced a number.
The present chapter produces numbers, and with them the first genuinely
rigid conventions of the book: once one decides what
$\int_\sigma\dprime x\wedge\ddprime x$ means on a unit segment,
every later sign---the sign in the Poincar\'e--Lelong formula, the sign of a
Chern current, the sign of the Dirichlet energy, the direction of the
inequality defining plurisubharmonicity---is determined.  A single misplaced
$(-1)^{n-1}$ here would propagate silently through
\Crefrange{chap:cld-forms}{chap:cld-examples} and would reverse the
positivity statements of \Cref{part:bfj}.  For this reason the chapter is
written with an unusual density of hand computations, and every convention it
fixes is recorded twice: once in the sign ledger of
\Cref{sec:ch11-ledger} and once in \Cref{app:normalizations}.

\section{The normalization problem}
\label{sec:ch11-normalization}

Ordinary Euclidean integration depends on a choice of Lebesgue measure.
Tropical integration has a canonical normalization because every cell comes
with a tangent lattice.  Three invariance tests govern the construction:
\begin{enumerate}
\item changing a lattice basis must not change the integral;
\item subdividing a cell must not change the integral;
\item pushing a cycle through an integer-affine map must introduce exactly
  the lattice index of \Cref{prop:tropical-pushforward-balanced}.
\end{enumerate}
The reader should keep these tests separate.  The first two concern one
calibrated support; the third compares different supports.

This chapter proves the polyhedral core of Stokes and Green formulae.  The
passage to the full CLD category is an imported analytic theorem, stated with
its hypotheses.  Compact support will always mean compact support in the
underlying polyhedral or analytic topology.

\subsection{What has to be normalized, and why}
\label{sec:ch11-what-normalized}

A top-degree superform on an $n$-dimensional real affine space is, at each
point, a real multiple of the Lagerberg orientation $\tau_n$ of
\eqref{eq:appJ-lagerberg-orientation}, and to integrate it one needs a
measure.  An ordinary top-degree differential form changes sign when the
orientation is reversed.  A superform of bidegree $(n,n)$ transforms by the
square of that determinant, so its sign is orientation independent, but a
volume normalization is still needed.  The tangent lattice supplies the
covolume-one measure.  Thus weighted lattices are indispensable data on a
tropical cycle, whereas an orientation is only a bookkeeping choice that
\Cref{lem:ch11-orientation-drops-out} shows to disappear from every formula
of the chapter.  In the Lagerberg calculus the canonical orientation of a
complex manifold is replaced by the canonical squaring, and the lattice
replaces the complex structure.

The third invariance test is of a different nature.  Tests $(1)$ and
$(2)$ say that the integral is attached to a calibrated support, so that
two combinatorial presentations of the same tropical cycle give the same
number.  Test $(3)$ compares two different supports through a map and is
therefore a compatibility, not an invariance; it is the tropical shadow of
the degree of a finite morphism, and it is the only one of the three in which
a lattice index appears.  Confusing $(2)$ and $(3)$ produces exactly the
multiplicity errors warned against in \Cref{chap:tropical-spaces}.

\subsection{Prerequisites, dependencies, and what is imported}
\label{sec:ch11-prerequisites}

The reader needs the superform calculus of \Cref{chap:superforms}---in
particular the differentials $\dprime$, $\ddprime$, the wedge sign
convention, and the constant $c_p=(-1)^{p(p-1)/2}$ of
\eqref{eq:appJ-cp}---and the polyhedral vocabulary of
\Cref{chap:tropical-spaces}: rational polyhedra, tangent lattices
$N_\sigma$, primitive normal generators $u_{\sigma/\tau}$, weights,
balancing, calibrations, and integer-affine pushforward.  From ordinary
analysis we use only Fubini's theorem, the fundamental theorem of calculus,
the affine change-of-variables formula, and smooth partitions of unity.  No
measure theory beyond Lebesgue measure on $\mathbf R^n$ is required, and no
tropical intersection theory beyond
\Cref{thm:tropical-stable-intersection} is assumed.

Everything in
\Crefrange{sec:ch11-normalization}{sec:ch11-tropical-intersection}
is proved.  Exactly one statement is imported:
\Cref{thm:cld-global-stokes}, the passage from a single calibrated
polyhedral complex to a general calibrated parlinear space, uses the
existence of calibrated chart covers with compatible refinements.  That is
the substantive analytic content of \cite[\S\S4.6 and 5.11]{CLD25} and we do
not reprove it.  The reader who is content with skeleta and tropical
charts---which is all that
\Crefrange{chap:cld-poincare-lelong}{chap:cld-examples} ever needs in
practice---may treat \Cref{thm:balanced-stokes} as the operative Stokes
theorem and read \Cref{thm:cld-global-stokes} as a statement about
generality.

\subsection{The plan of the chapter}
\label{sec:ch11-plan}

By the end of the chapter the reader should be able to
\begin{enumerate}
\item integrate a top-degree superform over a weighted rational polyhedral
  complex and check the answer against a lattice normalization;
\item write down the boundary trace of an $(n-1,n)$-form across a
  codimension-one face, with the correct sign, in any dimension;
\item recognize the balancing condition as the exact vanishing of a boundary
  defect vector, and hence as the exact hypothesis under which Stokes'
  formula has no interior terms;
\item derive the Green and Dirichlet formulae and match their signs against
  the metric-graph conventions of \Cref{sec:metric-graphs};
\item push a cycle forward along a proper integer-affine map and verify the
  projection formula by hand;
\item state precisely, and prove in the cases used later in this book, the
  comparison between the tropical intersection product of
  \Cref{thm:tropical-stable-intersection} and the product of currents.
\end{enumerate}
Each of these is tested on an explicit example, and
\Cref{sec:ch11-ledger} collects the outcomes.

\section{Top-degree integration}
\label{sec:ch11-top-degree}

Throughout the chapter $N$ is a free abelian group of finite rank $r$,
$N_{\mathbf R}=N\otimes_{\mathbf Z}\mathbf R$, and
$M=\operatorname{Hom}(N,\mathbf Z)$, as in \Cref{chap:tropical-spaces}.  A
weighted complex $(\mathscr C,m)$ is pure of dimension $n\le r$; its set of
facets is written $\mathscr C_n$ and its set of codimension-one cells
$\mathscr C_{n-1}$.  For a cell $\sigma$ we write $N_\sigma\subset N$ for its
saturated tangent lattice and $N_{\sigma,\mathbf R}$ for the corresponding
real subspace, so that the affine span of $\sigma$ is
$x_\sigma+N_{\sigma,\mathbf R}$ for any $x_\sigma\in\sigma$.

\subsection{The lattice measure of a cell}
\label{sec:ch11-lattice-measure}

\begin{definition}
\label{def:ch11-lattice-measure}
\index{lattice measure}
Let $\sigma$ be an $n$-dimensional rational polyhedron in $N_{\mathbf R}$.
The \emph{lattice measure} $\omega_\sigma$ of $\sigma$ is the unique
translation-invariant Borel measure on the affine span of $\sigma$ for which
the parallelotope spanned by one---equivalently any---lattice basis of
$N_\sigma$ has measure one.
\end{definition}

Existence and uniqueness are immediate: transport Lebesgue measure from
$\mathbf R^n$ by a lattice basis.  What has to be checked is that the
transported measure does not depend on the basis, and this is the first
of the three invariance tests.

\begin{lemma}
\label{lem:ch11-lattice-measure-welldefined}
The measure $\omega_\sigma$ is well defined, and it depends only on the
lattice $N_\sigma$, not on $\sigma$ itself nor on any orientation.  If
$\tau\subset\sigma$ is a face of dimension $n-1$, then
$\omega_\tau$ is the lattice measure of the saturated lattice $N_\tau$, and
for every $v\in N_\sigma$ whose class generates $N_\sigma/N_\tau$ one has,
in the affine span of $\sigma$,
\[
 \omega_\sigma
 =\int_{\mathbf R}\bigl(t_*\omega_\tau\bigr)\,\rd t,
\]
where $t_*\omega_\tau$ denotes the translate of $\omega_\tau$ by $tv$.  In
words: the lattice measure of $\sigma$ disintegrates over the lattice
measure of $\tau$ with unit transverse speed in the direction of a primitive
transverse generator.
\end{lemma}

\begin{proof}
Two lattice bases of $N_\sigma$ differ by a matrix in
$\operatorname{GL}(n,\mathbf Z)$, whose determinant is $\pm1$; the
transported Lebesgue measures therefore agree, because Lebesgue measure
transforms by $|\det|$.  The measure so obtained visibly depends only on
$N_\sigma$, and it is unsigned, so no orientation enters.

For the disintegration, note that $N_\tau$ is saturated in $N_\sigma$ by
\Cref{chap:tropical-spaces}, so the quotient $N_\sigma/N_\tau$ is free of
rank one.  Choose $v\in N_\sigma$ mapping to a generator; then a lattice
basis $f_1,\ldots,f_{n-1}$ of $N_\tau$ together with $v$ is a lattice basis
of $N_\sigma$.  In the resulting coordinates $(t,s_1,\ldots,s_{n-1})$ the
measure $\omega_\sigma$ is $\rd t\,\rd s$ and $\omega_\tau$ is $\rd s$ on
each slice $t=\text{const}$, which is the assertion.  Replacing $v$ by
$-v$ or by $v+w$ with $w\in N_\tau$ changes neither side.
\end{proof}

The disintegration formula is the only fact about lattice measures used in
the proof of Stokes' theorem, and it is worth isolating the reason it is
true: the transverse speed is set to one by \emph{primitivity} of the
transverse generator, not by any metric.  Using a non-primitive generator,
or a Euclidean unit normal, destroys the identity.  We return to this in
\Cref{ex:ch11-euclidean-normal-trap}.

\medskip

Let $\sigma\subset N_{\mathbf R}$ be an $n$-dimensional rational
polyhedron.  Choose a lattice basis $e_1,\ldots,e_n$ of $N_\sigma$ and dual
affine coordinates $x_1,\ldots,x_n$.

\begin{definition}
\label{def:ch11-polyhedral-integral}
Put
\[
\tau_n=\dprime x_1\wedge \ddprime x_1\wedge\cdots\wedge \dprime x_n\wedge \ddprime x_n.
\]
For a compactly supported $(n,n)$-form
\[
\alpha=f\tau_n
\]
on $\sigma$, set
\[
                           \int_\sigma\alpha
             =\int_\sigma f(x)\,\rd x_1\cdots \rd x_n.
\]
For a compactly supported $(n,n)$-form $\alpha$ on a pure
$n$-dimensional weighted complex $(\mathscr C,m)$, set
\[
 \int_{(\mathscr C,m)}\alpha
 =\sum_{\sigma\in\mathscr C_n}m_\sigma\int_\sigma\alpha.
\]
The sum is finite by local finiteness.
\end{definition}

If $\omega_\sigma$ denotes the covolume-one Lebesgue measure determined by
$N_\sigma$, the definition can be written invariantly as
\[
                    \int_\sigma f\tau_n
                    =\int_\sigma f\,\omega_\sigma.
\]
The interlaced form $\tau_n$ fixes the positivity sign.  Using the grouped
order of all primed factors followed by all double-primed factors would
multiply the positive generator by $(-1)^{n(n-1)/2}$ and would therefore
reverse positivity in dimensions two and three.

Before the invariance statements we record the two elementary identities on
which every sign in this chapter rests.  Both are consequences of the wedge
convention of \Cref{chap:superforms}, which places all primed factors before
all double-primed factors, and both are used constantly.

\begin{lemma}
\label{lem:ch11-interlaced-grouped}
Write $c_p=(-1)^{p(p-1)/2}$ as in \eqref{eq:appJ-cp}.  For any affine
coordinates $x_1,\ldots,x_n$,
\[
 \tau_n
 =c_n\,\dprime x_1\wedge\cdots\wedge\dprime x_n
      \wedge\ddprime x_1\wedge\cdots\wedge\ddprime x_n,
 \qquad
 \dprime x_1\wedge\cdots\wedge\ddprime x_n=c_n\tau_n .
\]
Moreover $c_nc_{n-1}=(-1)^{n-1}$, and the $2$-forms
$\dprime x_i\wedge\ddprime x_i$ commute with everything.
\end{lemma}

\begin{proof}
The first identity is the displayed count in \Cref{chap:superforms}: moving
$\ddprime x_i$ to the right past $\dprime x_{i+1},\ldots,\dprime x_n$ costs
$(-1)^{n-i}$, and $\sum_{i=1}^{n}(n-i)=n(n-1)/2$.  The second follows since
$c_n^2=1$.  For the third,
\[
 c_nc_{n-1}
 =(-1)^{\frac{n(n-1)}2+\frac{(n-1)(n-2)}2}
 =(-1)^{\frac{(n-1)(n+n-2)}2}
 =(-1)^{(n-1)^2}
 =(-1)^{n-1}.
\]
Finally $\dprime x_i\wedge\ddprime x_i$ has total degree two, hence is
central in the bigraded algebra by the commutation rule of
\Cref{chap:superforms}.
\end{proof}

The constant $c_nc_{n-1}=(-1)^{n-1}$ is the single most consequential
number in this chapter: it is exactly the sign appearing in the boundary
trace \eqref{eq:superform-boundary-traces}, and every sign error we have
encountered in drafts of this material was a lost copy of it.

\begin{lemma}
\label{lem:ch11-orientation-drops-out}
Let $x_1,\ldots,x_n$ and $y_1,\ldots,y_n$ be affine coordinates dual to two
lattice bases of $N_\sigma$, not assumed to induce the same orientation.
Then
\[
 \dprime x_1\wedge\ddprime x_1\wedge\cdots
   \wedge\dprime x_n\wedge\ddprime x_n
 =\dprime y_1\wedge\ddprime y_1\wedge\cdots
   \wedge\dprime y_n\wedge\ddprime y_n .
\]
Consequently $\tau_n$ depends only on the lattice $N_\sigma$, the integral
of \Cref{def:ch11-polyhedral-integral} is unsigned, and no orientation of
$\sigma$ enters $\int_\sigma\alpha$ for $\alpha\in A^{n,n}$.
\end{lemma}

\begin{proof}
Let $B\in\operatorname{GL}(n,\mathbf Z)$ be the matrix with
$y_i=\sum_jB_{ij}x_j$ up to constants.  By
\Cref{lem:ch11-interlaced-grouped} it suffices to compare the grouped
products.  Then
\[
 \dprime y_1\wedge\cdots\wedge\dprime y_n=\det(B)\,
 \dprime x_1\wedge\cdots\wedge\dprime x_n,
\]
and the same identity holds for the double-primed factors.  Hence the
grouped product is multiplied by $\det(B)^2=1$, since $\det(B)=\pm1$.  The
two interlaced products therefore agree.  The remaining assertions follow
because $\int_\sigma f\tau_n=\int_\sigma f\,\omega_\sigma$ and
$\omega_\sigma$ is unsigned by
\Cref{lem:ch11-lattice-measure-welldefined}.
\end{proof}

This is the precise sense in which orientation choices used in coordinate
descriptions are bookkeeping.  If one integrated an ordinary
$n$-form the sign $\det(B)=\pm1$ would survive and a global orientation of
the support would be indispensable.  Tropical cycles are generally not
orientable in that sense---the tropical line of
\Cref{fig:tropical-line-balancing} carries three rays meeting at a point and
admits no coherent orientation---so it is essential that the superform
calculus does not require one.  This is one of the structural reasons the
Lagerberg formalism, rather than ordinary differential forms, is the right
carrier for non-Archimedean curvature.

\begin{proposition}
\label{prop:integral-subdivision-independent}
Let $(\mathscr C,m)$ and $(\mathscr C',m')$ be pure $n$-dimensional weighted
rational polyhedral complexes representing the same tropical $n$-cycle, in
the sense of \Cref{chap:tropical-spaces}.  Then
$|\mathscr C|=|\mathscr C'|$ and, for every $(n,n)$-form $\alpha$ with
compact support on that set,
\[
 \int_{(\mathscr C,m)}\alpha=\int_{(\mathscr C',m')}\alpha .
\]
The common value depends only on the cycle and on the restriction of
$\alpha$ to the cells.
\end{proposition}

\begin{proof}
Basis independence is \Cref{lem:ch11-orientation-drops-out}: a change of
lattice basis has determinant $\pm1$, and the two labelled exterior blocks
contribute its square.

By definition of equivalence there is a weighted complex
$(\mathscr D,\ell)$ refining both, with each facet of $\mathscr D$ carrying
the weight of the unique facet of $\mathscr C$ (respectively of
$\mathscr C'$) containing it.  It therefore suffices to treat one refinement,
say $\mathscr D\to\mathscr C$.  Put
$K=\Supp\alpha\cap|\mathscr C|$.  Local finiteness and compactness imply
that only finitely many facets of either complex meet $K$.

Fix a facet $\sigma\in\mathscr C_n$ meeting $K$, and let
$\delta_1,\ldots,\delta_k$ be the facets of $\mathscr D$ contained in
$\sigma$ and meeting $K$.  By the subdivision lemma of
\Cref{chap:tropical-spaces}, their relative interiors are pairwise disjoint
and cover $\sigma$ wherever $\alpha$ is nonzero, up to a finite union of
cells of dimension less than $n$, hence up to an
$\omega_\sigma$-null set.  Each $\delta_j$ spans the same affine space as
$\sigma$, so $N_{\delta_j}=N_\sigma$ and therefore
$\omega_{\delta_j}=\omega_\sigma$; saturation makes the two tangent
lattices equal, rather than merely commensurable.  Writing
$\alpha=f\tau_n$ in coordinates adapted to $N_\sigma$, finite additivity
of the Lebesgue integral gives
\[
 \int_\sigma f\,\omega_\sigma
 =\sum_{j=1}^k\int_{\delta_j}f\,\omega_{\delta_j}.
\]
Multiplying by $m_\sigma=\ell_{\delta_j}$ and summing over
$\sigma\in\mathscr C_n$ yields the assertion.

Finally, the definition of $A^{n,n}$ on a polyhedral support
(\Cref{chap:tropical-spaces}) already identifies two ambient representatives
whose restrictions to all cells agree, so the number produced depends only on
those restrictions.
\end{proof}

\begin{warning}{Saturation is a hypothesis, not a convenience}
The proof used $N_{\delta}=N_\sigma$ for a subcell $\delta\subset\sigma$ of
the same dimension.  This is true because tangent lattices were defined as
saturations, $N_\sigma=N\cap\operatorname{span}_{\mathbf R}(\sigma-\sigma)$.
Had one defined $N_\sigma$ as the subgroup generated by integral vectors of
$\sigma-\sigma$, the two lattices could differ by a finite index and
subdivision would change the integral.  The same remark explains the
primitivity requirement on $u_{\sigma/\tau}$.
\end{warning}

\begin{example}[A non-unimodular parametrization]
\label{ex:ch11-nonunimodular}
Let $F:[0,1]\to[0,2]$ be $F(t)=2t$, with lattice $\mathbf Z$ on source and
target.  The pushforward cell has weight
$[\mathbf Z:2\mathbf Z]=2$.  For a compactly supported function $g$ on
$(0,2)$,
\[
 \int_{F_*[0,1]} g(x)\,\rd x
 =2\int_0^2g(x)\,\rd x,
\]
where the notation suppresses the two labelled top-degree factors.  For the
actual superform
$\alpha=g(x)\,\dprime x\wedge\ddprime x$, one has
\[
 F^*\alpha=4g(2t)\,\dprime t\wedge\ddprime t,
 \qquad
 \int_{[0,1]}F^*\alpha
 =2\int_0^2g(x)\,\rd x.
\]
The factor four comes from pulling back both labelled copies; the ordinary
change-of-variables factor removes one copy.  What remains is precisely the
lattice index in the pushforward weight.
\end{example}

\begin{example}[Positivity sign test]
\label{ex:ch11-positivity-sign}
With $c_n=(-1)^{n(n-1)/2}$,
\[
\tau_n=c_n\,\dprime x_1\wedge\cdots\wedge \dprime x_n\wedge
 \ddprime x_1\wedge\cdots\wedge \ddprime x_n.
\]
Thus a compactly supported strongly positive top form has nonnegative
integral.
\end{example}

The next example is the promised failure of the naive convention, in the
smallest dimension in which it can occur.  It should be worked through
completely; the analogous mistake at the level of boundary traces is
harder to see and is isolated in \Cref{ex:ch11-square-wrong-sign}.

\begin{example}[The grouped orientation is the wrong normalization]
\label{ex:ch11-grouped-wrong-sign}
Work on $\mathbf R^2$ with the standard lattice, and let
$\chi\ge0$ be a smooth compactly supported function.  Consider
\[
 \alpha=\chi\,\dprime x_1\wedge\ddprime x_1
       \wedge\dprime x_2\wedge\ddprime x_2
      =\chi\,\tau_2 ,
\]
which is strongly positive: it is
$\chi$ times the product of
$\dprime x_i\wedge J(\dprime x_i)$ for $i=1,2$, with a nonnegative smooth
coefficient.  Its integral is
$\int\chi\,\rd x_1\rd x_2\ge0$, as it must be.

Suppose instead that one had declared the \emph{grouped} monomial
\[
 \gamma=\dprime x_1\wedge\dprime x_2
       \wedge\ddprime x_1\wedge\ddprime x_2
\]
to be the positive top form and set
$\int(g\gamma)=\int g\,\rd x_1\rd x_2$.  By
\Cref{lem:ch11-interlaced-grouped}, $\gamma=c_2\tau_2=-\tau_2$, so this
convention assigns to the strongly positive form $\alpha$ the value
$-\int\chi\le0$.  Every positivity statement in the book---semipositivity of
a Chern current, nonnegativity of a Monge--Amp\`ere measure, the sign of the
Dirichlet form in \Cref{cor:tropical-dirichlet-sign}---would be reversed in
even dimension and preserved in odd dimension, which is precisely the sort of
dimension-dependent error that survives a spot check in $n=1$.  The
one-dimensional test $\tau_1=\gamma$ detects nothing.  The smallest honest
test is $n=2$.
\end{example}

\begin{example}[Euclidean normals are not lattice normals]
\label{ex:ch11-euclidean-normal-trap}
Let $N=\mathbf Z^2$ and let $\tau$ be the segment from $(1,0)$ to $(0,1)$.
Its tangent lattice is $N_\tau=\mathbf Z(1,-1)$, so $\omega_\tau$ assigns
mass one to $\tau$: the lattice length of $\tau$ is $1$, whereas its
Euclidean length is $\sqrt2$.  Let $\sigma$ be the triangle with vertices
$(0,0)$, $(1,0)$, $(0,1)$, so that $N_\sigma=\mathbf Z^2$ and
$\omega_\sigma(\sigma)=\tfrac12$.

The class of $(1,0)$ generates $N_\sigma/N_\tau$, and so does the class of
$(0,1)$, since $(1,0)-(0,1)=(1,-1)\in N_\tau$.  Either may be used as the
primitive transverse generator $v$ of
\Cref{lem:ch11-lattice-measure-welldefined}.  The Euclidean unit normal
$(1,1)/\sqrt2$ is not an integral vector at all, and the primitive integral
vector in that direction is $(1,1)$, which represents the class $2v$ in
$N_\sigma/N_\tau$ because $(1,1)=(1,0)+(0,1)$.  Using it would double every
boundary contribution across $\tau$.

The disintegration of \Cref{lem:ch11-lattice-measure-welldefined} can be
checked directly with $v=(1,0)$: the slices of $\sigma$ transverse to
$N_{\tau,\mathbf R}$ are the segments
$\sigma\cap\bigl(\langle\tau\rangle-tv\bigr)
 =\sigma\cap\{x_1+x_2=1-t\}$ for
$0\le t\le1$, measured inward from the hypotenuse, whose lattice lengths are
$1-t$, and
$\int_0^1(1-t)\,\rd t=\tfrac12=\omega_\sigma(\sigma)$.  Had one used
Euclidean lengths and the Euclidean unit normal the same computation would
read $\int_0^{1/\sqrt2}\sqrt2\,(1-\sqrt2 t)\,\rd t=\tfrac12$, which is again
correct---the two errors cancel---but the intermediate quantities are not
the ones appearing in the balancing condition, and the cancellation fails as
soon as two cells with different tangent lattices meet along $\tau$.  This is
why the primitive lattice normal, and not the Euclidean one, is the datum
carried by a calibration.
\end{example}

\begin{example}[A unimodular change of chart]
\label{ex:ch11-gl2z-test}
Let $C$ be the square $[0,1]^2\subset\mathbf R^2$ with weight $1$ and the
standard lattice, and let
$\alpha=\chi(x_1,x_2)\,\tau_2$.  Apply the integer-affine automorphism
$B(x_1,x_2)=(x_1+x_2,x_2)$, which lies in $\operatorname{SL}(2,\mathbf Z)$.
Then $B_*C$ is the parallelogram with vertices
$(0,0),(1,0),(2,1),(1,1)$, with weight one by
\Cref{prop:tropical-pushforward-balanced}, since $B$ is an isomorphism of
lattices.  Pulling back,
\[
 B^*(\dprime y_1\wedge\ddprime y_1\wedge\dprime y_2\wedge\ddprime y_2)
 =(\dprime x_1+\dprime x_2)\wedge(\ddprime x_1+\ddprime x_2)
   \wedge\dprime x_2\wedge\ddprime x_2
 =\tau_2 ,
\]
because the cross terms containing $\dprime x_2$ or $\ddprime x_2$ twice
vanish.  Hence
$\int_{B_*C}\alpha=\int_C B^*\alpha$, and both equal
$\int_{[0,1]^2}\chi(B(x))\,\rd x$.  The Euclidean areas of $C$ and $B(C)$
are both $1$; but if one repeats the computation with
$B'(x_1,x_2)=(2x_1,x_2)$, which is not unimodular, the image has Euclidean
area $2$ while the pushforward weight becomes
$[\mathbf Z^2:B'(\mathbf Z^2)]=2$ and the two effects again cancel in
\Cref{prop:ch11-pushforward-integral}.  Test $(1)$ and test $(3)$ of
\Cref{sec:ch11-normalization} are visible side by side here.
\end{example}

\section{Boundary normals and contraction}
\label{sec:ch11-boundary}

The divergence theorem in $\mathbf R^n$ has three ingredients: a boundary
measure, an outward normal, and a rule converting a vector field into a
flux.  In the superform calculus all three are packaged in a single
operation, contraction in one of the two labelled exterior blocks.  We
therefore begin by fixing that operation, because its sign convention is not
the classical one for ordinary forms and the difference is exactly the
source of the constant $(-1)^{n-1}$.

\subsection{Contraction in a labelled block}
\label{sec:ch11-contraction}

\begin{definition}
\label{def:ch11-contractions}
\index{contraction!in a labelled block}
Let $v\in V$ be a vector.  On a pure tensor
$\eta\otimes\xi\in\bigwedge^pV^*\otimes\bigwedge^qV^*$ set
\[
 \iota'_v(\eta\otimes\xi)=(\iota_v\eta)\otimes\xi,
 \qquad
 \iota''_v(\eta\otimes\xi)=\eta\otimes(\iota_v\xi),
\]
where $\iota_v$ is the ordinary interior product, and extend by linearity
and by $C^\infty$-linearity in the coefficients.  Thus $\iota'_v$ lowers the
first degree and $\iota''_v$ the second, and \emph{no sign is incurred when
$\iota''_v$ passes the primed block}.
\end{definition}

The convention that $\iota''_v$ ignores the primed factors is the one that
makes the traces below uniform in the dimension; the alternative convention,
in which $\iota''_v$ is treated as an odd operator on the total degree and
acquires $(-1)^p$, differs from ours by exactly $(-1)^p$ and would move that
sign into \eqref{eq:superform-boundary-traces}.  Both are used in the
literature.  Ours is fixed for the whole book: it is the convention of
\Cref{def:ch09-contraction}, whose graded variant
$\widetilde\iota''_v=(-1)^p\iota''_v$ is the alternative just described
(\Cref{rem:ch09-contraction-convention}), and it is recorded in
\Cref{sec:appJ-currents}.

\begin{lemma}
\label{lem:ch11-contraction-basic}
The operators $\iota'_v$ and $\iota''_v$ have the following properties.
\begin{enumerate}
\item\label{itm:ch11-contr:linear}
  They are $\mathbf R$-linear in $v$, and
  $\iota'_v\iota'_v=\iota''_v\iota''_v=0$ and
  $\iota'_v\iota''_w=\iota''_w\iota'_v$ for all $v,w$.
\item\label{itm:ch11-contr:coord}
  In affine coordinates, for
  $\alpha=\sum_{I,J}\alpha_{IJ}\,\dprime x_I\wedge\ddprime x_J$,
  \[
   \iota'_v\alpha
   =\sum_{I,J}\alpha_{IJ}\,\bigl(\iota_v\dprime x_I\bigr)
     \wedge\ddprime x_J,
   \qquad
   \iota''_v\alpha
   =\sum_{I,J}\alpha_{IJ}\,\dprime x_I\wedge
     \bigl(\iota_v\ddprime x_J\bigr),
  \]
  where $\iota_v$ denotes the ordinary interior product in the labelled
  exterior algebra concerned.  In particular each of $\iota'_v$ and
  $\iota''_v$ is an antiderivation \emph{within its own block}: if
  $\lambda_1,\ldots,\lambda_k$ are covectors, then
  \[
   \iota''_v\bigl(\ddprime\lambda_1\wedge\cdots\wedge\ddprime\lambda_k\bigr)
   =\sum_{a=1}^k(-1)^{a-1}\langle\lambda_a,v\rangle\,
     \ddprime\lambda_1\wedge\cdots\widehat{\ddprime\lambda_a}\cdots
     \wedge\ddprime\lambda_k .
  \]
  \item\label{itm:ch11-contr:leibniz}
  For $\alpha\in A^{p,q}$ and $\beta\in A^{p',q'}$,
  \[
   \begin{aligned}
   \iota'_v(\alpha\wedge\beta)
     &=(\iota'_v\alpha)\wedge\beta
       +(-1)^{p+q}\alpha\wedge(\iota'_v\beta),\\
   \iota''_v(\alpha\wedge\beta)
     &=(-1)^{p'}(\iota''_v\alpha)\wedge\beta
       +(-1)^{q}\alpha\wedge(\iota''_v\beta).
   \end{aligned}
  \]
\end{enumerate}
\end{lemma}

\begin{proof}
\cref{itm:ch11-contr:linear} and \cref{itm:ch11-contr:coord} are immediate
from \Cref{def:ch11-contractions} and the corresponding properties of the
ordinary interior product, applied in one tensor factor at a time;
commutation holds because the two operators act on different factors.

\cref{itm:ch11-contr:leibniz}
Write $\alpha=\eta\otimes\xi$ and $\beta=\eta'\otimes\xi'$ with
$\deg\eta=p$, $\deg\xi=q$, $\deg\eta'=p'$, $\deg\xi'=q'$, and recall from
\Cref{chap:superforms} that
$(\eta\otimes\xi)\wedge(\eta'\otimes\xi')
=(-1)^{qp'}(\eta\wedge\eta')\otimes(\xi\wedge\xi')$.  Applying $\iota'_v$ and
the ordinary antiderivation rule in the first factor,
\[
 \iota'_v(\alpha\wedge\beta)
 =(-1)^{qp'}\bigl[(\iota_v\eta)\wedge\eta'
   +(-1)^p\eta\wedge(\iota_v\eta')\bigr]\otimes(\xi\wedge\xi').
\]
Now $(\iota'_v\alpha)\wedge\beta=(-1)^{qp'}((\iota_v\eta)\wedge\eta')
\otimes(\xi\wedge\xi')$ and
$\alpha\wedge(\iota'_v\beta)=(-1)^{q(p'-1)}(\eta\wedge(\iota_v\eta'))
\otimes(\xi\wedge\xi')$, so the second term equals
$(-1)^{qp'+p-q(p'-1)}\alpha\wedge(\iota'_v\beta)
=(-1)^{p+q}\alpha\wedge(\iota'_v\beta)$.

For $\iota''_v$ the same computation in the second factor gives
\[
 \iota''_v(\alpha\wedge\beta)
 =(-1)^{qp'}(\eta\wedge\eta')\otimes
   \bigl[(\iota_v\xi)\wedge\xi'+(-1)^q\xi\wedge(\iota_v\xi')\bigr],
\]
while $(\iota''_v\alpha)\wedge\beta=(-1)^{(q-1)p'}(\eta\wedge\eta')
\otimes((\iota_v\xi)\wedge\xi')$ and
$\alpha\wedge(\iota''_v\beta)=(-1)^{qp'}(\eta\wedge\eta')
\otimes(\xi\wedge(\iota_v\xi'))$.  Comparing coefficients gives the
factors $(-1)^{qp'-(q-1)p'}=(-1)^{p'}$ and $(-1)^{q}$.
\end{proof}

The asymmetry between the two Leibniz rules in
\cref{itm:ch11-contr:leibniz} is a consequence of the convention that
$\iota''_v$ ignores the primed block: an operator acting on the second
factor of a tensor product cannot be a graded derivation for the total
degree unless it is given a traversal sign.  Only
\cref{itm:ch11-contr:linear} and \cref{itm:ch11-contr:coord} are used in the
proofs below, so nothing depends on the asymmetry; it is recorded because a
reader who reconstructs the boundary traces with the other convention will
meet the transposed formulas.

\subsection{The outward normal and the boundary traces}
\label{sec:ch11-traces}

Let $\tau$ be a codimension-one face of $\sigma$. Choose a lift
$n_{\sigma,\tau}\in N_\sigma$ of the outward primitive normal class, the
negative of the inward class $u_{\sigma/\tau}$.  No orientation of either
cell is needed: the traces below are top-degree superforms on $\tau$, and
their two labelled determinant factors make their integrals orientation
independent.
Write $\iota'_v$ and $\iota''_v$ for contraction in the primed and
double-primed exterior factors.  The boundary traces compatible with the
interlaced top-degree form are
\begin{equation}
\label{eq:superform-boundary-traces}
\begin{aligned}
 \operatorname{tr}'_{\sigma,\tau}(\eta)
   &=(-1)^{n-1}\left.(\iota''_{n_{\sigma,\tau}}\eta)\right|_\tau,
       &&\eta\in A^{n-1,n},\\
 \operatorname{tr}''_{\sigma,\tau}(\theta)
   &=-\left.(\iota'_{n_{\sigma,\tau}}\theta)\right|_\tau,
       &&\theta\in A^{n,n-1}.
\end{aligned}
\end{equation}
The labelled factor is forced by bidegree: a form of bidegree $(n-1,n)$
has no spare primed factor to contract.  The traces do not depend on the
chosen lift.  Indeed, changing it by a vector tangent to $\tau$ gives zero
after restriction, because the contracted exterior block had degree $n$.

We now prove the three properties of \eqref{eq:superform-boundary-traces}
that will be used without further comment: they are well defined, they see
the form only through its restriction to $\sigma$, and they are computed in
adapted coordinates by a formula with no residual sign.

\begin{lemma}
\label{lem:ch11-trace-welldefined}
Let $\tau$ be a codimension-one face of the $n$-dimensional rational
polyhedron $\sigma$, and let $\eta\in A^{n-1,n}$, $\theta\in A^{n,n-1}$ be
superforms defined on a neighbourhood of $\tau$ in $N_{\mathbf R}$.  Then:
\begin{enumerate}
\item\label{itm:ch11-trace:lift}
  $\operatorname{tr}'_{\sigma,\tau}(\eta)$ and
  $\operatorname{tr}''_{\sigma,\tau}(\theta)$ do not depend on the choice of
  the lift $n_{\sigma,\tau}\in N_\sigma$ of the outward primitive class in
  $N_\sigma/N_\tau$;
\item\label{itm:ch11-trace:restriction}
  they vanish if $\eta$, respectively $\theta$, restricts to zero on
  $\sigma$; hence they descend to the space of superforms on the polyhedral
  support in the sense of \Cref{chap:tropical-spaces};
\item\label{itm:ch11-trace:target}
  they are $(n-1,n-1)$-forms on $\tau$ and therefore may be integrated
  against $\omega_\tau$ with no further choice.
\end{enumerate}
\end{lemma}

\begin{proof}
\cref{itm:ch11-trace:lift}
Two lifts differ by a vector $w\in N_\tau$.  By linearity of $\iota''$ in
the vector, the difference of the two candidate traces is
$(-1)^{n-1}(\iota''_w\eta)|_\tau$.  Evaluate this $(n-1,n-1)$-form on
tangent vectors of $\tau$: the result is $\eta$ evaluated on $n-1$ primed
vectors tangent to $\tau$ and on the $n$ double-primed vectors
$w,v_1,\ldots,v_{n-1}$, all lying in the space $N_{\tau,\mathbf R}$ of
dimension $n-1$.  Such an $n$-tuple is linearly dependent, so the
alternating form vanishes.  The same argument applies to
$\iota'_w\theta$, where the primed block now has degree $n$.

\cref{itm:ch11-trace:restriction}
Suppose $\eta|_\sigma=0$, that is, $\eta$ vanishes whenever all its
arguments are tangent to $\sigma$.  Evaluating
$(\iota''_{n_{\sigma,\tau}}\eta)|_\tau$ on tangent vectors of $\tau$
inserts the vectors $v_1,\ldots,v_{n-1}\in N_{\tau,\mathbf R}$ in the primed
block and $n_{\sigma,\tau},w_1,\ldots,w_{n-1}$ in the double-primed block.
All of these lie in $N_{\sigma,\mathbf R}$, because the lift was chosen in
$N_\sigma$ and $N_\tau\subset N_\sigma$.  Hence the value is zero.  Again
the argument for $\theta$ is identical.  This is exactly the point at which
the requirement $n_{\sigma,\tau}\in N_\sigma$, rather than merely
$n_{\sigma,\tau}\in N$, is used.

\cref{itm:ch11-trace:target}
Contraction lowers one degree, so $\iota''_{n_{\sigma,\tau}}\eta$ has
bidegree $(n-1,n-1)$ and restricting to $\tau$ preserves the bidegree.  Since
$\dim\tau=n-1$, this is a top-degree form on $\tau$ and
\Cref{def:ch11-lattice-measure} applies.
\end{proof}

\begin{lemma}
\label{lem:ch11-trace-normal-form}
Let $e_1,\ldots,e_n$ be a lattice basis of $N_\sigma$ such that
$e_2,\ldots,e_n$ is a lattice basis of $N_\tau$ and $e_1$ represents the
outward primitive class, and let $x_1,\ldots,x_n$ be the dual affine
coordinates, so that $\sigma$ lies in $\{x_1\le0\}$ near $\tau$ and
$\tau\subset\{x_1=0\}$.  Write
\[
 \eta=c_n\sum_{i=1}^n(-1)^{i-1}f_i\,
   \dprime x_1\wedge\cdots\widehat{\dprime x_i}\cdots\wedge\dprime x_n
   \wedge\ddprime x_1\wedge\cdots\wedge\ddprime x_n ,
\]
\[
 \theta=c_n\sum_{i=1}^n(-1)^{n+i-1}g_i\,
   \dprime x_1\wedge\cdots\wedge\dprime x_n
   \wedge\ddprime x_1\wedge\cdots\widehat{\ddprime x_i}
        \cdots\wedge\ddprime x_n .
\]
Then, with $\tau_{n-1}$ formed from $x_2,\ldots,x_n$,
\[
 \dprime\eta=\Bigl(\sum_i\frac{\partial f_i}{\partial x_i}\Bigr)\tau_n,
 \qquad
 \ddprime\theta=\Bigl(\sum_i\frac{\partial g_i}{\partial x_i}\Bigr)\tau_n,
\]
and
\[
 \operatorname{tr}'_{\sigma,\tau}(\eta)=f_1|_\tau\,\tau_{n-1},
 \qquad
 \operatorname{tr}''_{\sigma,\tau}(\theta)=g_1|_\tau\,\tau_{n-1}.
\]
\end{lemma}

\begin{proof}
Every $(n-1,n)$-form on the span of $\sigma$ can be written as displayed,
because the monomials appearing are a basis of
$\bigwedge^{n-1}\oplus\bigwedge^n$ in these coordinates; the factor
$c_n(-1)^{i-1}$ is a normalization chosen to make the conclusion clean.

\emph{The differentials.}  By the definition of $\dprime$,
\[
 \dprime\eta
 =c_n\sum_{i,j}(-1)^{i-1}\frac{\partial f_i}{\partial x_j}\,
   \dprime x_j\wedge\dprime x_1\wedge\cdots\widehat{\dprime x_i}\cdots
   \wedge\dprime x_n\wedge\ddprime x_1\wedge\cdots\wedge\ddprime x_n .
\]
Only $j=i$ survives, and
$\dprime x_i\wedge\dprime x_1\wedge\cdots\widehat{\dprime x_i}\cdots
\wedge\dprime x_n=(-1)^{i-1}\dprime x_1\wedge\cdots\wedge\dprime x_n$, so
the two signs $(-1)^{i-1}$ cancel and
\[
 \dprime\eta=c_n\Bigl(\sum_i\frac{\partial f_i}{\partial x_i}\Bigr)
  \dprime x_1\wedge\cdots\wedge\dprime x_n
   \wedge\ddprime x_1\wedge\cdots\wedge\ddprime x_n
  =\Bigl(\sum_i\frac{\partial f_i}{\partial x_i}\Bigr)\tau_n
\]
by \Cref{lem:ch11-interlaced-grouped} and $c_n^2=1$.  For $\theta$, the
defining sign of $\ddprime$ on a form whose primed degree is $n$ contributes
$(-1)^n$; only $j=i$ survives; and
$\ddprime x_i\wedge\ddprime x_1\wedge\cdots\widehat{\ddprime x_i}\cdots
\wedge\ddprime x_n=(-1)^{i-1}\ddprime x_1\wedge\cdots\wedge\ddprime x_n$.
The total sign is
$(-1)^n(-1)^{n+i-1}(-1)^{i-1}=(-1)^{2n+2i-2}=1$, whence the second identity.

\emph{The traces.}  Apply $\iota''_{e_1}$ to $\eta$.  By
\Cref{def:ch11-contractions} the primed block is untouched and
$\iota_{e_1}(\ddprime x_1\wedge\cdots\wedge\ddprime x_n)
=\ddprime x_2\wedge\cdots\wedge\ddprime x_n$, so
\[
 \iota''_{e_1}\eta
 =c_n\sum_i(-1)^{i-1}f_i\,
   \dprime x_1\wedge\cdots\widehat{\dprime x_i}\cdots\wedge\dprime x_n
   \wedge\ddprime x_2\wedge\cdots\wedge\ddprime x_n .
\]
Restricting to $\tau$ annihilates $\dprime x_1$, hence kills every summand
with $i\ne1$.  There remains
\[
 \begin{aligned}
 (\iota''_{e_1}\eta)|_\tau
 &=c_nf_1|_\tau\,
   \dprime x_2\wedge\cdots\wedge\dprime x_n
   \wedge\ddprime x_2\wedge\cdots\wedge\ddprime x_n\\
 &=c_nc_{n-1}f_1|_\tau\,\tau_{n-1}
  =(-1)^{n-1}f_1|_\tau\,\tau_{n-1},
 \end{aligned}
\]
using \Cref{lem:ch11-interlaced-grouped} twice.  Multiplying by the
prefactor $(-1)^{n-1}$ of \eqref{eq:superform-boundary-traces} gives
$\operatorname{tr}'_{\sigma,\tau}(\eta)=f_1|_\tau\,\tau_{n-1}$.

For $\theta$, apply $\iota'_{e_1}$; now the primed block loses $\dprime x_1$
and the restriction to $\tau$ kills every summand containing
$\ddprime x_1$, that is, every $i\ne1$.  There remains
  \[
   \begin{aligned}
   (\iota'_{e_1}\theta)|_\tau
   & =c_n(-1)^{n}g_1|_\tau\,
     \dprime x_2\wedge\cdots\wedge\dprime x_n
     \wedge\ddprime x_2\wedge\cdots\wedge\ddprime x_n\\
   & =(-1)^{n-1}(-1)^{n}g_1|_\tau\,\tau_{n-1}
    =-g_1|_\tau\,\tau_{n-1},
   \end{aligned}
  \]
and the prefactor $-1$ in \eqref{eq:superform-boundary-traces} produces
$\operatorname{tr}''_{\sigma,\tau}(\theta)=g_1|_\tau\,\tau_{n-1}$.
\end{proof}

Two consequences deserve to be stated separately, because they are what make
the formula usable.

\begin{corollary}
\label{cor:ch11-trace-flux}
Let $\sigma$ be an $n$-dimensional rational polyhedron with tangent lattice
 $N_\sigma$.  There are canonical isomorphisms of $C^\infty$-modules
 \[
   \begin{aligned}
   A^{n-1,n}\bigl(N_{\sigma,\mathbf R}\bigr)
   &\;\xrightarrow{\ \sim\ }\;C^\infty\otimes N_{\sigma,\mathbf R},
   &\eta&\longmapsto F_\eta,\\
   A^{n,n-1}\bigl(N_{\sigma,\mathbf R}\bigr)
   &\;\xrightarrow{\ \sim\ }\;C^\infty\otimes N_{\sigma,\mathbf R},
   &\theta&\longmapsto G_\theta,
   \end{aligned}
 \]
depending only on the lattice $N_\sigma$ and characterized in the adapted
coordinates of \Cref{lem:ch11-trace-normal-form} by
$F_\eta=\sum_if_ie_i$ and $G_\theta=\sum_ig_ie_i$.  With this notation
\[
 \dprime\eta=(\operatorname{div}F_\eta)\,\tau_n,
 \qquad
 \ddprime\theta=(\operatorname{div}G_\theta)\,\tau_n,
\]
and, for a codimension-one face $\tau\prec\sigma$,
\[
 \operatorname{tr}'_{\sigma,\tau}(\eta)=\lambda_\eta\,\tau_{n-1},
 \qquad
 \operatorname{tr}''_{\sigma,\tau}(\theta)=\lambda_\theta\,\tau_{n-1},
\]
where $\lambda_\eta,\lambda_\theta\in C^\infty(\tau)$ are the functions
determined by
\[
 F_\eta\equiv\lambda_\eta\,n_{\sigma,\tau},
 \qquad
 G_\theta\equiv\lambda_\theta\,n_{\sigma,\tau}
 \qquad\text{in }N_{\sigma,\mathbf R}/N_{\tau,\mathbf R}.
\]
In particular the two traces obey the \emph{same} rule, and both measure
outward flux with respect to the primitive lattice normal.
\end{corollary}

\begin{proof}
Canonically $A^{n-1,n}(V)=\bigwedge^{n-1}V^*\otimes\bigwedge^nV^*$ and
$\bigwedge^{n-1}V^*\cong V\otimes\bigwedge^nV^*$, so
$A^{n-1,n}(V)\cong V\otimes(\bigwedge^nV^*)^{\otimes2}$.  A lattice
$N_\sigma$ trivializes $(\bigwedge^nV^*)^{\otimes2}$ canonically, since two
lattice bases differ by a determinant $\pm1$ whose square is $1$; this is
\Cref{lem:ch11-orientation-drops-out} again, and it produces the first
isomorphism.  Exchanging the roles of the two exterior factors gives the
second.  In adapted coordinates the trivializing element is $\tau_n$, and the
normal forms displayed in \Cref{lem:ch11-trace-normal-form} are exactly the
images of $\sum_if_ie_i$ and $\sum_ig_ie_i$.

The formulas for $\dprime\eta$ and $\ddprime\theta$ were proved in
\Cref{lem:ch11-trace-normal-form}.  For the traces, choose adapted
coordinates with $e_1=n_{\sigma,\tau}$ and $e_2,\ldots,e_n$ a lattice basis
of $N_\tau$.  Then $F_\eta\equiv f_1e_1$ modulo $N_{\tau,\mathbf R}$, so
$\lambda_\eta=f_1|_\tau$, and \Cref{lem:ch11-trace-normal-form} gives
$\operatorname{tr}'_{\sigma,\tau}(\eta)=f_1|_\tau\,\tau_{n-1}$.  Both
$\lambda_\eta$ and the trace were defined without reference to the adapted
basis, so the identity holds in general.  The argument for $\theta$ is
identical.
\end{proof}

\begin{warning}{The flux coefficient is a quotient, not an inner product}
The function $\lambda_\eta$ is \emph{not} the Euclidean inner product of
$F_\eta$ with a unit normal.  It is the coefficient of the class of $F_\eta$
in the rank-one real vector space $N_{\sigma,\mathbf R}/N_{\tau,\mathbf R}$
with respect to the basis furnished by the primitive outward class.  In
particular, replacing $n_{\sigma,\tau}$ by twice itself halves
$\lambda_\eta$, and no metric intervenes.  On the hypotenuse of the standard
triangle of \Cref{ex:ch11-euclidean-normal-trap} the two prescriptions differ
by a factor $\sqrt2$.
\end{warning}

\begin{corollary}
\label{cor:ch11-trace-additive}
Fix $\tau$ and a superform $\eta\in A^{n-1,n}$ defined near $\tau$.  The
assignment
\[
 v\longmapsto (\iota''_v\eta)|_\tau
\]
is a linear map from $N_{\mathbf R}$ to the one-dimensional space of
$(n-1,n-1)$-forms at a point of $\tau$, and it kills
$N_{\tau,\mathbf R}$.  It therefore factors through
$N_{\mathbf R}/N_{\tau,\mathbf R}$.
\end{corollary}

\begin{proof}
Linearity is \Cref{lem:ch11-contraction-basic}, and the vanishing on
$N_{\tau,\mathbf R}$ was proved in
\Cref{lem:ch11-trace-welldefined}\,\cref{itm:ch11-trace:lift}, whose
argument did not use $v\in N_\sigma$.  The target is one dimensional because
$\bigwedge^{n-1}N_\tau^*\otimes\bigwedge^{n-1}N_\tau^*$ is.
\end{proof}

\Cref{cor:ch11-trace-additive} is the technical heart of the balancing
argument in \Cref{sec:ch11-balancing}: it converts a sum of boundary
contributions over the facets containing $\tau$ into a single contraction
against the sum of their weighted normal vectors.

\begin{lemma}
\label{lem:ch11-segment-stokes}
For $\sigma=[a,b]$ with its standard lattice coordinate $x$ and
$\eta=f\,\ddprime x$,
\[
                       \int_\sigma \dprime\eta=f(b)-f(a).
\]
The outward normals at $b$ and $a$ are $+1$ and $-1$ respectively, and the
right-hand side is the sum of the traces in
\eqref{eq:superform-boundary-traces}.
\end{lemma}
\begin{proof}
Since $\dprime\eta=f'\,\dprime x\wedge \ddprime x$, the assertion is the fundamental theorem
of calculus. This fixes all later boundary signs.
\end{proof}

It is worth carrying out the verification of the last sentence in complete
detail once, since it is the base case of every induction below.  Take
$n=1$, so $c_1=1$ and $\tau_1=\dprime x\wedge\ddprime x$, and write
$\eta=f\,\ddprime x$, which is the normal form of
\Cref{lem:ch11-trace-normal-form} with $f_1=f$.  The face $\tau=\{b\}$ has
$N_\tau=0$ and $\tau_0=1$, and the outward primitive normal is
$n_{\sigma,b}=+1$; hence
\[
 \operatorname{tr}'_{\sigma,\{b\}}(\eta)
 =(-1)^{1-1}\,\iota''_{+1}\bigl(f\,\ddprime x\bigr)\big|_{\{b\}}
 =f(b).
\]
At $\tau=\{a\}$ the outward normal is $-1$ and the same computation gives
$-f(a)$.  The sum is $f(b)-f(a)$, which is $\int_a^bf'\,\rd x$.  The
lattice measure of a point is the counting measure of mass one, so
$\int_{\{b\}}f(b)\tau_0=f(b)$; no length enters.

\begin{example}[A segment with a weight]
\label{ex:ch11-weighted-segment}
Let $\Gamma$ be the interval $[0,\ell]\subset\mathbf R$ with the standard
lattice $\mathbf Z$ and weight $m\in\mathbf Z\setminus\{0\}$.  For
$\eta=f\,\ddprime x$ one gets
\[
 \int_{(\Gamma,m)}\dprime\eta=m\bigl(f(\ell)-f(0)\bigr).
\]
Two features of the general theory are already present.  First, the answer is
insensitive to the Euclidean length $\ell$: only the values at the two
endpoints appear, because $\dprime\eta$ is an exact top form.  Second, the
weight multiplies the boundary terms as well as the interior integral;
this is what allows weights of opposite sign to cancel across a shared
endpoint, which is the one-dimensional case of balancing.  Concretely, if
$\Gamma_1=[-\ell,0]$ and $\Gamma_2=[0,\ell]$ carry weights $m_1$ and $m_2$
and $f$ vanishes near $\pm\ell$, then
\[
 \int_{(\Gamma_1,m_1)}\dprime\eta+\int_{(\Gamma_2,m_2)}\dprime\eta
 =m_1f(0)-m_2f(0)=(m_1-m_2)f(0),
\]
which vanishes for all $f$ precisely when $m_1=m_2$.  The primitive outward
normals at $0$ are $+1$ from $\Gamma_1$ and $-1$ from $\Gamma_2$, so the
condition $m_1-m_2=0$ is literally $\sum_\sigma m_\sigma n_{\sigma,\tau}=0$:
the balancing condition of \Cref{chap:tropical-spaces} at the vertex $0$.
\end{example}

\begin{example}[Boundary signs on a rectangle]
\label{ex:ch11-rectangle-boundary}
On $\sigma=[0,a]\times[0,b]$ with the standard lattice coordinates, take
\[
 \eta=-f(x,y)\,\dprime y\wedge\ddprime x\wedge\ddprime y.
\]
Only the $x$-derivative contributes to
$\dprime\eta$.  The two vertical faces have outward normals
$-\partial_x$ and $+\partial_x$, so their boundary integrals occur with
opposite signs.  The horizontal faces contribute zero because contraction
in the double-primed block with $\partial_y$ leaves the \emph{primed} factor
$\dprime y$, which vanishes on the horizontal tangent line; equivalently the
flux vector $F_\eta=(f,0)$ of \Cref{cor:ch11-trace-flux} is tangent to those
faces.  The resulting
identity is the fundamental theorem of calculus in $x$, integrated over
$y$.
\end{example}

The next example carries out the same computation with explicit numbers and
then repeats it with the prefactor $(-1)^{n-1}$ removed.  It is the smallest
complete demonstration that the prefactor is not cosmetic.

\begin{example}[The naive trace has the wrong sign in even dimension]
\label{ex:ch11-square-wrong-sign}
Let $\sigma=[0,1]^2\subset\mathbf R^2$ with $N=\mathbf Z^2$, and take
\[
 \eta=-x_1\,\dprime x_2\wedge\ddprime x_1\wedge\ddprime x_2
 \ \in A^{1,2}.
\]
This is the normal form of \Cref{lem:ch11-trace-normal-form} with $n=2$,
$c_2=-1$, $f_1=x_1$ and $f_2=0$; the associated vector field is
$F=(x_1,0)$, of divergence $1$.  Accordingly
\[
 \dprime\eta=\tau_2,
 \qquad
 \int_\sigma\dprime\eta=\omega_\sigma(\sigma)=1 .
\]

Now compute the four boundary traces from
\eqref{eq:superform-boundary-traces}, using
\Cref{cor:ch11-trace-flux}.  The outward primitive normals are
$-e_1,\;+e_1,\;-e_2,\;+e_2$ on the faces
$\{x_1=0\}$, $\{x_1=1\}$, $\{x_2=0\}$, $\{x_2=1\}$, and each face has
lattice length one.  On $\{x_2=1\}$ the tangent lattice is
$\mathbf Ze_1$ and $F=x_1e_1\equiv0$ modulo it, so $\lambda_\eta=0$; the
same holds on $\{x_2=0\}$, and the two horizontal faces contribute nothing.
On $\{x_1=1\}$ the tangent lattice is $\mathbf Ze_2$, the outward primitive
class is that of $e_1$, and $F=x_1e_1\equiv1\cdot e_1$, so
$\lambda_\eta=1$.  On $\{x_1=0\}$ the outward class is that of $-e_1$ and
$F=x_1e_1=0$ there, so $\lambda_\eta=0$.  The total is
\[
 \sum_\tau\int_\tau\operatorname{tr}'_{\sigma,\tau}(\eta)=1+0+0+0=1
 =\int_\sigma\dprime\eta ,
\]
as \Cref{thm:polyhedral-stokes} asserts.

Suppose now that one had defined the boundary trace by the naive recipe
\[
 \widetilde{\operatorname{tr}}'_{\sigma,\tau}(\eta)
 =\left.\bigl(\iota''_{n_{\sigma,\tau}}\eta\bigr)\right|_\tau,
\]
that is, without the prefactor $(-1)^{n-1}$.  On $\{x_1=1\}$ the
computation in the proof of \Cref{lem:ch11-trace-normal-form} gives
$(\iota''_{e_1}\eta)|_\tau=(-1)^{n-1}f_1\tau_{n-1}=-\tau_1$, whence
\[
 \sum_\tau\int_\tau\widetilde{\operatorname{tr}}'_{\sigma,\tau}(\eta)
 =-1\neq1 .
\]
Stokes' formula would fail by a sign.  In dimension one the two recipes
agree, because $(-1)^{n-1}=1$; in dimension three they agree again.  The
naive convention is therefore consistent with every check performed on
segments, on curves, and on threefolds, and fails exactly on surfaces---the
case in which most of the interesting examples of this book live.  We
record the moral: \emph{a boundary convention must be tested in even
dimension}.

The same example also separates two other frequently confused choices.  If
one replaces the outward normal $n_{\sigma,\tau}$ by the inward normal
$u_{\sigma/\tau}$, every trace changes sign and the total becomes $-1$; if
one replaces the primitive normal by a Euclidean unit normal on a face
whose tangent lattice is not generated by a Euclidean unit vector, the
individual boundary contributions are rescaled by irrational factors and no
longer combine into the balancing sum, as in
\Cref{ex:ch11-euclidean-normal-trap}.  Of these three possible errors, only
the second is visible to a dimension-one test.
\end{example}

\begin{theorem}\label{thm:polyhedral-stokes}
Let $\sigma$ be a rational $n$-polytope. For every smooth
$(n-1,n)$-form $\eta$ defined near $\sigma$,
\[
                         \int_\sigma \dprime\eta
          =\sum_{\substack{\tau\prec\sigma\\
                           \operatorname{codim}_\sigma\tau=1}}\int_\tau
             \operatorname{tr}'_{\sigma,\tau}(\eta).
\]
For $\theta\in A^{n,n-1}$ the analogous formula holds with $\ddprime$ and
$\operatorname{tr}''_{\sigma,\tau}(\theta)$.
\end{theorem}
\begin{proof}
By linearity it suffices to write
\[
\eta=c_n\sum_{i=1}^n(-1)^{i-1}f_i\,
\dprime x_1\wedge\cdots\widehat{\dprime x_i}\cdots\wedge \dprime x_n
\wedge \ddprime x_1\wedge\cdots\wedge \ddprime x_n.
\]
Then
\[
\dprime\eta=\left(\sum_i\frac{\partial f_i}{\partial x_i}\right)\tau_n.
\]
The asserted formula is exactly the divergence theorem for the vector field
$(f_1,\ldots,f_n)$ with lattice-normalized measure.  The proof has four
steps: the classical Euclidean divergence theorem, the conversion of
Euclidean boundary data into lattice boundary data, the identification of
the resulting expression with the trace
\eqref{eq:superform-boundary-traces}, and the double-primed variant.  Throughout, $x_1,\ldots,x_n$ are the
affine coordinates dual to a lattice basis $e_1,\ldots,e_n$ of $N_\sigma$
fixed once and for all, so that $\omega_\sigma$ is Lebesgue measure
$\rd x=\rd x_1\cdots\rd x_n$; we equip $N_{\sigma,\mathbf R}$ with the
auxiliary Euclidean structure making $e_1,\ldots,e_n$ orthonormal.  This
Euclidean structure is a device: it disappears from the final formula, as
Step 2 shows.

\smallskip
\emph{Step 1: the Euclidean divergence theorem.}
Let $F=(f_1,\ldots,f_n)$, let $\tau$ run over the codimension-one faces of
$\sigma$, let $\nu_\tau$ be the Euclidean outward unit normal along $\tau$
and $\mathcal H^{n-1}$ the Euclidean surface measure.  Then
\begin{equation}
\label{eq:ch11-euclidean-divergence}
 \int_\sigma\operatorname{div}F\,\rd x
 =\sum_\tau\int_\tau\langle F,\nu_\tau\rangle\,\rd\mathcal H^{n-1}.
\end{equation}
This is the classical Gauss--Green theorem for a bounded open set with
Lipschitz boundary, applied to the interior of $\sigma$; see
\cite[\S5.8]{EvansGariepy}.

We include the reduction to one variable.  On the standard simplex
\[
 \Delta_n=\{x_i\geq0,\ \textstyle\sum_i x_i\leq1\},
\]
Fubini and the fundamental theorem of calculus give, for each $i$,
\[
\begin{split}
 \int_{\Delta_n}\frac{\partial f_i}{\partial x_i}\,\rd x
 =\int_{\Delta_{n-1}}
 \bigl[
 f_i(x_1,\ldots,1-\!\sum_{j\ne i}x_j,\ldots,x_n)\\
 {}-f_i(x_1,\ldots,0,\ldots,x_n)
 \bigr]\,\rd x_{\widehat{\imath}} ,
\end{split}
\]
where $\Delta_{n-1}$ is the image of $\Delta_n$ in the hyperplane
$\{x_i=0\}$; the interchange of the order of integration is legitimate
because $f_i$ is smooth near the compact set $\Delta_n$.  The second term is
the flux across the coordinate face, whose Euclidean outward unit normal is
$-e_i$ and whose surface measure is $\rd x_{\widehat\imath}$.  On the
remaining face $\{\sum_jx_j=1\}$ one has
$\nu=n^{-1/2}(1,\ldots,1)$, and the projection $p_i$ forgetting the $i$th
coordinate is a bijection onto $\Delta_{n-1}$ with
$(p_i)_*\mathcal H^{n-1}=n^{1/2}\,\rd x_{\widehat\imath}$, because the angle
between $\nu$ and $e_i$ has cosine $n^{-1/2}$.  Hence the $i$th first term is
$\int n^{-1/2}f_i\,\rd\mathcal H^{n-1}$, and summing over $i$ gives
$\int n^{-1/2}\bigl(\sum_if_i\bigr)\rd\mathcal H^{n-1}
=\int\langle F,\nu\rangle\,\rd\mathcal H^{n-1}$.  Thus
\eqref{eq:ch11-euclidean-divergence} holds on $\Delta_n$; the general
polytope is covered by the cited theorem, and also, for a reader who
prefers an elementary route, by triangulating and observing that interior
faces occur twice with opposite Euclidean outward unit normals and equal
surface measures, so that their contributions cancel.

\smallskip
\emph{Step 2: from Euclidean to lattice boundary data.}
Fix a codimension-one face $\tau$ and let $d_\tau>0$ be the covolume of the
lattice $N_\tau$ inside the hyperplane $N_{\tau,\mathbf R}$ for the auxiliary
Euclidean structure.  By definition of the lattice measure,
\[
 \rd\mathcal H^{n-1}\big|_\tau=d_\tau\,\omega_\tau .
\]
Next, $N_\tau\oplus\mathbf Z\,n_{\sigma,\tau}=N_\sigma$, because
$n_{\sigma,\tau}$ represents a generator of the rank-one group
$N_\sigma/N_\tau$.  Computing the covolume of $N_\sigma$ in two ways,
\[
 1=\operatorname{covol}(N_\sigma)
  =\operatorname{covol}(N_\tau)\cdot
   \bigl|\langle n_{\sigma,\tau},\nu_\tau\rangle\bigr|
  =d_\tau\,\bigl|\langle n_{\sigma,\tau},\nu_\tau\rangle\bigr|,
\]
the middle equality being the base-times-height formula for the
parallelotope spanned by a lattice basis of $N_\tau$ together with
$n_{\sigma,\tau}$.  As $n_{\sigma,\tau}$ points outward, the pairing is
positive and $\langle n_{\sigma,\tau},\nu_\tau\rangle=d_\tau^{-1}$.

Let $\lambda_\eta$ be the flux coefficient of
\Cref{cor:ch11-trace-flux}, so that $F\equiv\lambda_\eta\,n_{\sigma,\tau}$
modulo $N_{\tau,\mathbf R}$.  Since $\nu_\tau$ annihilates
$N_{\tau,\mathbf R}$,
\[
 \langle F,\nu_\tau\rangle
 =\lambda_\eta\,\langle n_{\sigma,\tau},\nu_\tau\rangle
 =\lambda_\eta\,d_\tau^{-1},
\]
whence
\[
 \int_\tau\langle F,\nu_\tau\rangle\,\rd\mathcal H^{n-1}
 =\int_\tau\lambda_\eta\,d_\tau^{-1}\cdot d_\tau\,\omega_\tau
 =\int_\tau\lambda_\eta\,\omega_\tau .
\]
The auxiliary Euclidean structure has cancelled, as it had to: the left-hand
side of \eqref{eq:ch11-euclidean-divergence} does not depend on it.

\smallskip
\emph{Step 3: the sign.}
The factor $(-1)^{n-1}$ in the first trace is the quotient
$c_n/c_{n-1}$ between the interlaced top forms in dimensions $n$ and
$n-1$.  This shows directly that the first trace in
\eqref{eq:superform-boundary-traces} gives the displayed flux terms.
Precisely, \Cref{cor:ch11-trace-flux} gives
$\operatorname{tr}'_{\sigma,\tau}(\eta)=\lambda_\eta\,\tau_{n-1}$, and hence
$\int_\tau\operatorname{tr}'_{\sigma,\tau}(\eta)
=\int_\tau\lambda_\eta\,\omega_\tau$, which by Step 2 is the flux term of
\eqref{eq:ch11-euclidean-divergence}.  Combining Steps 1--3 and
$\int_\sigma\dprime\eta=\int_\sigma\operatorname{div}F\,\rd x$ proves the
$\dprime$ formula.  This is the only step in which a sign is created; Steps
1 and 2 involve positive factors only.

\smallskip
\emph{Step 4: the double-primed formula.}
For the other differential write
\[
 \theta=c_n\sum_{i=1}^n(-1)^{n+i-1}g_i\,
 \dprime x_1\wedge\cdots\wedge\dprime x_n
 \wedge\ddprime x_1\wedge\cdots
       \widehat{\ddprime x_i}\cdots\wedge\ddprime x_n.
\]
The defining sign in $\ddprime$ gives
\[
 \ddprime\theta=
 \left(\sum_i\frac{\partial g_i}{\partial x_i}\right)\tau_n.
\]
On a face with outward normal $e_i$, primed contraction followed by
restriction has coefficient $-g_i$ relative to the interlaced boundary
form; hence the minus sign in
$\operatorname{tr}''_{\sigma,\tau}$ produces the outward flux $g_i$.
Both assertions are \Cref{lem:ch11-trace-normal-form}, so
$\operatorname{tr}''_{\sigma,\tau}(\theta)=\lambda_\theta\,\tau_{n-1}$ with
$G_\theta\equiv\lambda_\theta n_{\sigma,\tau}$, and Steps 1--2 apply
verbatim with $F$ replaced by $G_\theta$.  Equivalently, apply the
involution $J$ of \Cref{chap:superforms}: it exchanges $\dprime$ and
$\ddprime$, carries $A^{n,n-1}$ to $A^{n-1,n}$, fixes $\tau_n$ up to the
sign $(-1)^n$ and $\tau_{n-1}$ up to $(-1)^{n-1}$, and conjugates $\iota'_v$
into the \emph{graded} contraction $\widetilde\iota''_v=(-1)^p\iota''_v$ of
\Cref{def:ch09-contraction}, not into $\iota''_v$
(\Cref{lem:ch09-contraction-basic}\,(3)); the two normalizations in
\eqref{eq:superform-boundary-traces} differ by exactly the resulting
factor.
\end{proof}

\begin{remark}[Why the interior cancellation of Step~1 is not yet balancing]
\label{rem:ch11-step4-vs-balancing}
In the triangulation argument of Step~1 the two simplices adjacent to an
interior face carry the
\emph{same} tangent lattice and the \emph{same} weight, namely the weight of
$\sigma$, so the cancellation is automatic.  On a general weighted complex
neither is true: adjacent facets may have different weights, and more than
two facets may meet along a codimension-one cell.  The cancellation then
becomes a hypothesis, and that hypothesis is balancing.  This is the
distinction which \Cref{sec:ch11-balancing} makes precise, and it is the
reason the present theorem is stated for a single polytope.
\end{remark}

\begin{example}[A triangle computation]
\label{ex:ch11-triangle}
Let
\[
 \Delta=\{(x,y)\in\mathbf R_{\geq0}^2\mid x+y\leq1\}
\]
with the standard lattice.  For the vector field
$F=(x,y)$, the divergence is $2$, so
$\int_\Delta\operatorname{div}(F)\,\rd x\,\rd y=1$.
The flux across the coordinate faces is zero.  On $x+y=1$, the primitive
outward normal in the quotient lattice is represented by $(1,0)$, or equally
by $(0,1)$ modulo the tangent vector $(1,-1)$, and the lattice-normalized
boundary parameter has length one.  The contraction prescribed above gives
total flux one.  This illustrates why one uses a primitive quotient normal,
not the Euclidean unit normal $(1,1)/\sqrt2$.
\end{example}

The arithmetic of the last example repays a second look, because it is the
first place where lattice and Euclidean normalizations visibly disagree and
where their disagreement nevertheless cancels.  With $F=(x,y)$ one has
$\operatorname{div}F=2$ and $\omega_\Delta(\Delta)=\tfrac12$, so the interior
integral is $1$.  On the hypotenuse the lattice measure gives mass $1$
(\Cref{ex:ch11-euclidean-normal-trap}), the outward primitive class is that
of $e_1$, and on the hypotenuse $F=xe_1+ye_2\equiv(x+y)e_1=1\cdot e_1$, so
$\lambda_F=1$ and the flux is $1$.  In the Euclidean bookkeeping the
hypotenuse has length $\sqrt2$, the unit normal is $(1,1)/\sqrt2$, and the
flux density is $(x+y)/\sqrt2=1/\sqrt2$, giving $\sqrt2\cdot(1/\sqrt2)=1$
again.  The two computations agree here because the ratio of Euclidean to
lattice measure on the face is exactly the ratio of the primitive to the
Euclidean unit normal; that is a general identity for a \emph{single}
polytope.  It fails as soon as two cells whose tangent lattices sit
differently inside $N$ meet along $\tau$, because then only the lattice
version is a sum of comparable terms; see
\Cref{ex:ch11-tripod-stokes}.

\begin{corollary}
\label{cor:polyhedral-integration-by-parts}
Let $\sigma$ be a polytope and suppose that one of the smooth functions
$f,g$ has support in the relative interior of $\sigma$. Then
\[
 \int_\sigma f\,\dprimeddprime g\wedge\alpha
 =-\int_\sigma \dprime f\wedge\ddprime g\wedge\alpha
\]
for every $\dprime$-closed form $\alpha$ of complementary bidegree for which
the products have top degree.
\end{corollary}
\begin{proof}
Apply \Cref{thm:polyhedral-stokes} to
$f\,\ddprime g\wedge\alpha$.  The support hypothesis kills every boundary
term, and $\dprime\alpha=0$ kills the term obtained by differentiating
$\alpha$.
\end{proof}

\section{Balancing and global cancellation}
\label{sec:ch11-balancing}

We have proved a divergence theorem on one polyhedron.  A tropical cycle is
many polyhedra glued along their faces, each carrying a weight, and the
boundary terms of adjacent facets must cancel if the total is to be zero.
The purpose of this section is to isolate the exact obstruction to that
cancellation and to observe that it is the balancing condition of
\Cref{chap:tropical-spaces}---not a sufficient condition for cancellation,
nor a convenient one, but precisely the necessary and sufficient one.  This
is the conceptual payoff of the chapter: balancing is not an axiom imposed
because tropicalizations happen to satisfy it, but the unique hypothesis
under which the calculus of superforms has a Stokes theorem.

Before proving anything we extend the previous theorem to unbounded cells,
which is needed because tropical cycles are rarely compact.

\begin{corollary}
\label{cor:ch11-stokes-unbounded}
Let $\sigma\subset N_{\mathbf R}$ be an $n$-dimensional rational polyhedron,
not assumed bounded, and let $\eta\in A^{n-1,n}$ be smooth on a
neighbourhood of $\sigma$ with $\sigma\cap\Supp\eta$ compact.  Then
\[
 \int_\sigma\dprime\eta
 =\sum_{\substack{\tau\prec\sigma\\ \operatorname{codim}_\sigma\tau=1}}
   \int_\tau\operatorname{tr}'_{\sigma,\tau}(\eta),
\]
all but finitely many terms being zero, and similarly for $\ddprime$.
\end{corollary}

\begin{proof}
Choose a rational box $B$, that is a product of bounded rational intervals
in coordinates dual to a lattice basis of $N$, whose interior contains the
compact set $\sigma\cap\Supp\eta$.  Then $\sigma\cap B$ is a rational
polytope with $N_{\sigma\cap B}=N_\sigma$, and
\Cref{thm:polyhedral-stokes} applies to it.  Its codimension-one faces are
of two kinds: those contained in a codimension-one face $\tau$ of $\sigma$,
and those contained in $\partial B$.  A face of the second kind is disjoint
from $\Supp\eta$, because $\Supp\eta\cap\sigma$ lies in the interior of
$B$, so its trace vanishes.  A face $\rho$ of the first kind satisfies
$N_\rho=N_\tau$ and lies on the same side of the affine span of $\tau$ as
$\sigma$ does, so it carries the same outward primitive normal; grouping the
$\rho$'s inside each $\tau$ and using additivity of $\omega_\tau$ replaces
$\sum_\rho\int_\rho$ by $\sum_\tau\int_\tau$, since $\eta$ vanishes on
$\tau\setminus B$.  Finally $\int_{\sigma\cap B}\dprime\eta
=\int_\sigma\dprime\eta$ for the same support reason.  Local finiteness of
the face poset of $\sigma$ gives the finiteness assertion.
\end{proof}

\subsection{The boundary defect of a weighted complex}
\label{sec:ch11-defect}

\begin{definition}
\label{def:ch11-defect}
\index{boundary defect}\index{balancing!as vanishing of a defect}
Let $(\mathscr C,m)$ be a pure $n$-dimensional weighted rational polyhedral
complex in $N_{\mathbf R}$ and let $\tau\in\mathscr C_{n-1}$.  The
\emph{boundary defect} of $(\mathscr C,m)$ along $\tau$ is
\[
 \mathbf N(\tau)
 =\sum_{\substack{\sigma\in\mathscr C_n\\ \sigma\supset\tau}}
    m_\sigma\,n_{\sigma,\tau}
 \;=\;-\sum_{\substack{\sigma\in\mathscr C_n\\ \sigma\supset\tau}}
    m_\sigma\,u_{\sigma/\tau}
 \qquad\in N/N_\tau .
\]
\end{definition}

The defect is well defined because each $n_{\sigma,\tau}=-u_{\sigma/\tau}$
is a canonical element of $N_\sigma/N_\tau\subset N/N_\tau$, and the sum is
finite by local finiteness of $\mathscr C$.  Comparing with
\Cref{chap:tropical-spaces}, the complex is balanced if and only if
$\mathbf N(\tau)=0$ for every $\tau\in\mathscr C_{n-1}$; the sign of the
defect is immaterial for that statement, but it matters for the formula
below, and we have fixed it so that $\mathbf N$ is a sum of \emph{outward}
normals, matching the outward convention of the boundary traces.

This is exactly Chambert-Loir and Ducros' \emph{discordance} of a calibration
along a codimension-one face, and their \emph{harmonious} spaces are those
of vanishing defect; see \cite[\S\S4.6.7--4.6.12]{CLD25}.  The vocabularies
differ because their calibration is defined on an abstract parlinear space,
whereas we work with an embedded complex; the content is the same.

\begin{warning}{Two independent sign choices}
Balancing is usually written with inward primitive generators,
$\sum_\sigma m_\sigma u_{\sigma/\tau}=0$; boundary traces are written with
outward normals.  Since the two families differ by a global factor $-1$,
the \emph{condition} balanced is insensitive to the choice, but the
\emph{defect} changes sign.  Whenever a residual boundary term is retained,
as in \Cref{thm:ch11-defect-stokes} and in the metric-graph comparison of
\Cref{thm:ch11-graph-green}, one must say which convention is in force.  Ours
is outward, always.
\end{warning}

\begin{theorem}
\label{thm:ch11-defect-stokes}
Let $(\mathscr C,m)$ be a pure $n$-dimensional weighted rational polyhedral
complex in $N_{\mathbf R}$, not assumed balanced.  For every
$\eta\in A^{n-1,n}_c(|\mathscr C|)$,
\begin{equation}
\label{eq:ch11-defect-stokes}
 \int_{(\mathscr C,m)}\dprime\eta
 =(-1)^{n-1}\sum_{\tau\in\mathscr C_{n-1}}
   \int_\tau\bigl(\iota''_{\mathbf N(\tau)}\eta\bigr)\big|_\tau ,
\end{equation}
where $\iota''_{\mathbf N(\tau)}$ means contraction against any lift of
$\mathbf N(\tau)$ to $N$, and only finitely many terms are nonzero.  The
analogous formula holds for $\theta\in A^{n,n-1}_c$ with $\ddprime$,
$\iota'$, and the global prefactor $-1$.
\end{theorem}

\begin{proof}
Since $\Supp\eta$ is compact and $\mathscr C$ is locally finite, only
finitely many cells meet $\Supp\eta$, so every sum below is finite and no
convergence question arises.

Apply \Cref{cor:ch11-stokes-unbounded} to each facet $\sigma$ and multiply
by $m_\sigma$:
\[
 \int_{(\mathscr C,m)}\dprime\eta
 =\sum_{\sigma\in\mathscr C_n}m_\sigma
   \sum_{\substack{\tau\prec\sigma\\ \operatorname{codim}=1}}
   \int_\tau\operatorname{tr}'_{\sigma,\tau}(\eta).
\]
Every codimension-one face of a facet of $\mathscr C$ is again a cell of
$\mathscr C$, because a polyhedral complex is closed under passage to faces.
We may therefore exchange the order of summation and group by
$\tau\in\mathscr C_{n-1}$:
\[
 \int_{(\mathscr C,m)}\dprime\eta
 =\sum_{\tau\in\mathscr C_{n-1}}\ \int_\tau
  \sum_{\sigma\supset\tau}m_\sigma
   \operatorname{tr}'_{\sigma,\tau}(\eta).
\]
Fix $\tau$.  By definition
$\operatorname{tr}'_{\sigma,\tau}(\eta)
=(-1)^{n-1}(\iota''_{n_{\sigma,\tau}}\eta)|_\tau$, and by
\Cref{cor:ch11-trace-additive} the assignment
$v\mapsto(\iota''_v\eta)|_\tau$ is linear and factors through
$N_{\mathbf R}/N_{\tau,\mathbf R}$.  Hence
\[
 \sum_{\sigma\supset\tau}m_\sigma\operatorname{tr}'_{\sigma,\tau}(\eta)
 =(-1)^{n-1}\Bigl(\iota''_{\sum_\sigma m_\sigma n_{\sigma,\tau}}
    \eta\Bigr)\Big|_\tau
 =(-1)^{n-1}\bigl(\iota''_{\mathbf N(\tau)}\eta\bigr)\big|_\tau ,
\]
which is \eqref{eq:ch11-defect-stokes}.  The argument for $\ddprime$ is
identical, with $\operatorname{tr}''_{\sigma,\tau}(\theta)
=-(\iota'_{n_{\sigma,\tau}}\theta)|_\tau$ and with
\Cref{cor:ch11-trace-additive} applied to $\iota'$, whose proof used only
that the contracted block has degree $n$.
\end{proof}

Formula \eqref{eq:ch11-defect-stokes} is the whole mechanism in one line.
The integral of an exact top form over a weighted complex is a sum of
codimension-one integrals, one for each codimension-one cell, each of which
is the contraction of $\eta$ against a single vector: the weighted sum of
outward normals.  Nothing else survives.  In particular the answer does not
depend on how the facets are arranged around $\tau$, only on that one
vector.

\begin{theorem}\label{thm:balanced-stokes}
Let $(\mathscr C,m)$ be a pure $n$-dimensional balanced rational polyhedral
complex. If $\eta\in A^{n-1,n}_c(|\mathscr C|)$, then
\[
                         \int_{(\mathscr C,m)}\dprime\eta=0.
\]
Similarly $\int \ddprime\theta=0$ for
$\theta\in A^{n,n-1}_c$.
\end{theorem}
\begin{proof}
In \eqref{eq:ch11-defect-stokes} every defect $\mathbf N(\tau)$ vanishes by
balancing.  The double-primed formula follows from its analogous defect
identity.
\end{proof}

In the language of \Cref{thm:ch11-defect-stokes} the proof is one line: the
right-hand side of \eqref{eq:ch11-defect-stokes} is a sum of terms
$\iota''_{\mathbf N(\tau)}\eta$, and $\mathbf N(\tau)=0$.  The following
converse shows that nothing weaker will do.

\begin{theorem}
\label{thm:ch11-balancing-iff-stokes}
Let $(\mathscr C,m)$ be a pure $n$-dimensional weighted rational polyhedral
complex in $N_{\mathbf R}$.  The following are equivalent.
\begin{enumerate}
\item\label{itm:ch11-bal:balanced}
  $(\mathscr C,m)$ is balanced.
\item\label{itm:ch11-bal:dprime}
  $\displaystyle\int_{(\mathscr C,m)}\dprime\eta=0$ for every
  $\eta\in A^{n-1,n}_c(|\mathscr C|)$.
\item\label{itm:ch11-bal:ddprime}
  $\displaystyle\int_{(\mathscr C,m)}\ddprime\theta=0$ for every
  $\theta\in A^{n,n-1}_c(|\mathscr C|)$.
\item\label{itm:ch11-bal:local}
  For every $\tau\in\mathscr C_{n-1}$ there is a neighbourhood $U$ of a
  point of the relative interior of $\tau$ in $N_{\mathbf R}$ such that
  $\int_{(\mathscr C,m)}\dprime\eta=0$ for every $\eta$ which is the
  restriction to $|\mathscr C|$ of a smooth $(n-1,n)$-form with compact
  support contained in $U$.
\end{enumerate}
In particular it suffices, in \cref{itm:ch11-bal:dprime} and
\cref{itm:ch11-bal:ddprime}, to test against restrictions of compactly
supported ambient forms.
\end{theorem}

\begin{proof}
\cref{itm:ch11-bal:balanced}$\Rightarrow$\cref{itm:ch11-bal:dprime} and
\cref{itm:ch11-bal:balanced}$\Rightarrow$\cref{itm:ch11-bal:ddprime} are
\Cref{thm:balanced-stokes}, and
\cref{itm:ch11-bal:dprime}$\Rightarrow$\cref{itm:ch11-bal:local} is trivial.
It remains to prove
\cref{itm:ch11-bal:local}$\Rightarrow$\cref{itm:ch11-bal:balanced}; the
implication \cref{itm:ch11-bal:ddprime}$\Rightarrow$\cref{itm:ch11-bal:balanced}
then follows by the same argument with $\iota'$ in place of $\iota''$, or by
applying the involution $J$.

Fix $\tau\in\mathscr C_{n-1}$ and a point $x_0$ in its relative interior.

\emph{Step 1: an adapted neighbourhood.}
We claim that $x_0$ has a neighbourhood $U$ in $N_{\mathbf R}$ meeting no
cell of $\mathscr C$ except $\tau$ and the facets containing $\tau$.
Indeed, let $\rho$ be a cell of $\mathscr C$ with $x_0\in\rho$.  Then
$\rho\cap\tau$ is a common face of $\rho$ and $\tau$ containing $x_0$;
since $x_0$ lies in the relative interior of $\tau$, the only face of $\tau$
containing $x_0$ is $\tau$ itself, so $\tau\subset\rho$.  Thus the cells
through $x_0$ are exactly $\tau$ and the cells containing it, and by purity
the ones of dimension $n$ are the facets $\sigma\supset\tau$.  The remaining
cells form a locally finite family of closed sets avoiding $x_0$, so some
neighbourhood $U$ misses all of them.  Shrinking $U$, we may assume
$\overline U$ compact.

\emph{Step 2: the local Stokes identity.}
Let $\eta$ be supported in $U$.  By \Cref{thm:ch11-defect-stokes} only
$\tau$ contributes, so
\[
 0=\int_{(\mathscr C,m)}\dprime\eta
  =(-1)^{n-1}\int_\tau\bigl(\iota''_{\mathbf N(\tau)}\eta\bigr)\big|_\tau .
\]

\emph{Step 3: a test form detecting the defect.}
Suppose $\mathbf N(\tau)\neq0$ in $N/N_\tau$, and let
$w\in N$ be a lift, so that $w\notin N_{\tau,\mathbf R}$ because $N_\tau$ is
saturated.  Choose $\xi\in M_{\mathbf R}$ with
$\xi|_{N_{\tau,\mathbf R}}=0$ and $\langle\xi,w\rangle=1$; this is possible
because $w$ is not in the subspace $N_{\tau,\mathbf R}$.  Choose also
$\xi_2,\ldots,\xi_n\in M_{\mathbf R}$ restricting to the basis of
$N_\tau^*$ dual to a lattice basis of $N_\tau$.  Finally let
$\chi\ge0$ be smooth, supported in $U$, with $\chi(x_0)>0$, and set
\[
 \eta=\chi\;
   \dprime\xi_2\wedge\cdots\wedge\dprime\xi_n
   \wedge\ddprime\xi\wedge\ddprime\xi_2\wedge\cdots\wedge\ddprime\xi_n
 \;\in A^{n-1,n}_c .
\]
Contracting in the double-primed block by
\Cref{lem:ch11-contraction-basic}\,\cref{itm:ch11-contr:coord},
\[
 \iota''_w\eta
 =\chi\,\langle\xi,w\rangle\,
   \dprime\xi_2\wedge\cdots\wedge\dprime\xi_n
   \wedge\ddprime\xi_2\wedge\cdots\wedge\ddprime\xi_n
  \;+\;\chi\,\dprime\xi_2\wedge\cdots\wedge\dprime\xi_n\wedge
     \ddprime\xi\wedge R,
\]
where $R$ collects the remaining $n-2$ double-primed factors with signs;
every term of the second group still contains the factor $\ddprime\xi$.  Restricting to $\tau$ annihilates that factor, because
$\xi$ vanishes on $N_{\tau,\mathbf R}$.  Hence, by
\Cref{lem:ch11-interlaced-grouped},
\[
 \bigl(\iota''_w\eta\bigr)\big|_\tau
 =\chi\,c_{n-1}\,\tau_{n-1}
 \quad\text{and}\quad
 \int_\tau\bigl(\iota''_w\eta\bigr)\big|_\tau
 =c_{n-1}\int_\tau\chi\,\omega_\tau\neq0,
\]
the last integral being strictly positive since $\chi\ge0$ is positive at
the interior point $x_0$ of $\tau$.  This contradicts Step 2.  Therefore
$\mathbf N(\tau)=0$, and as $\tau$ was arbitrary the complex is balanced.
\end{proof}

\begin{remark}[What the converse really says]
\label{rem:ch11-converse-meaning}
\Cref{thm:ch11-balancing-iff-stokes} promotes balancing from a combinatorial
axiom to a characterization.  Three consequences are worth recording.

First, balancing is a \emph{local} condition detected by test forms
concentrated near one codimension-one cell; this is implication
\cref{itm:ch11-bal:local}$\Rightarrow$\cref{itm:ch11-bal:balanced}.  Hence one
cannot repair an unbalanced complex by global cancellation between distant
cells.

Second, the equivalence is insensitive to the weights being integral: the
proof of Step 3 only used that $\mathbf N(\tau)$ is a real vector, so the
same statement holds for $\mathbf R$-weighted complexes.  Integrality of
weights is needed elsewhere---for lattice indices in pushforward, and for
the interpretation of weights as multiplicities of analytic
cycles---but not here.

Third, a codimension-one cell contained in a single facet has
$\mathbf N(\tau)=m_\sigma n_{\sigma,\tau}\neq0$, since $m_\sigma\neq0$ and
$n_{\sigma,\tau}$ is primitive.  Thus a balanced complex automatically has
no boundary in the naive sense; accordingly,
\Cref{thm:balanced-stokes} needs no separate boundary hypothesis.
\end{remark}

\subsection{The cancellation computed by hand}
\label{sec:ch11-hand-cancellation}

The proofs above are short because \Cref{cor:ch11-trace-additive} does the
work.  It is nevertheless worth seeing the cancellation happen in
coordinates at least once, in the smallest interesting case: a
one-dimensional complex.  There $n=1$, so $\tau$ is a point, $N_\tau=0$,
$\tau_0=1$, $\omega_\tau$ is the unit mass at that point, and
$\mathbf N(\tau)\in N$ is an honest vector.

\begin{example}[Stokes on a segment]
\label{ex:ch11-segment-stokes}
Let $C$ be the interval $[0,\ell]$ in $N_{\mathbf R}=\mathbf R$, with
$N=\mathbf Z$ and weight $m$.  Let $\eta=f\,\ddprime x$ be smooth.  Then
$\dprime\eta=f'\tau_1$ and
\[
 \int_{(C,m)}\dprime\eta=m\int_0^\ell f'(x)\,\rd x=m\bigl(f(\ell)-f(0)\bigr).
\]
The complex has two codimension-one cells, $\{0\}$ and $\{\ell\}$, with
defects $\mathbf N(\{0\})=-m$ and $\mathbf N(\{\ell\})=+m$.  Formula
\eqref{eq:ch11-defect-stokes} with $n=1$ reads
\[
 \int_{(C,m)}\dprime\eta
 =\iota''_{-m}\eta\big|_{0}+\iota''_{+m}\eta\big|_{\ell}
 =-mf(0)+mf(\ell),
\]
in agreement.  Neither defect vanishes: a segment is not balanced, and
Stokes' formula on it retains genuine boundary terms.  This is the correct
state of affairs; the reader should resist the temptation to call a segment
a tropical cycle.
\end{example}

\begin{example}[Stokes on a tripod, and the cancellation]
\label{ex:ch11-tripod-stokes}
Let $N=\mathbf Z^2$ and let $C$ be the tropical line of
\Cref{fig:tropical-line-balancing} truncated at lattice distance $L$: three
segments
\[
 \sigma_j=\{t\,u_j: 0\le t\le L\},
 \qquad
 u_1=(1,0),\quad u_2=(0,1),\quad u_3=(-1,-1),
\]
each of weight $1$, meeting at the vertex $\tau=\{0\}$, together with the
three endpoints $p_j=Lu_j$.  Each $u_j$ is primitive, so
$N_{\sigma_j}=\mathbf Zu_j$ and the lattice coordinate on $\sigma_j$ is the
parameter $t$.

Let $\eta=a_1\,\ddprime x_1+a_2\,\ddprime x_2$ with $a_i$ smooth on
$\mathbf R^2$; every element of $A^{0,1}$ near $C$ is of this form.  On
$\sigma_j$, restriction to the tangent line gives
$\eta|_{\sigma_j}=\langle a,u_j\rangle\,\ddprime t$, where
$\langle a,u_j\rangle=\sum_ia_iu_{j,i}$, so
\[
 \dprime\eta\big|_{\sigma_j}
 =\frac{\rd}{\rd t}\langle a(tu_j),u_j\rangle\,\dprime t\wedge\ddprime t,
 \qquad
 \int_{\sigma_j}\dprime\eta
 =\langle a(p_j),u_j\rangle-\langle a(0),u_j\rangle .
\]
Summing over $j$,
\begin{equation}
\label{eq:ch11-tripod-sum}
 \int_C\dprime\eta
 =\sum_{j=1}^3\langle a(p_j),u_j\rangle
  -\Bigl\langle a(0),\sum_{j=1}^3u_j\Bigr\rangle .
\end{equation}
The three primitive directions sum to zero, so the interior term
disappears and only the three endpoint terms remain.  This is the promised
cancellation, computed without any machinery: it happened because a single
vector, $\sum_ju_j$, is zero.

Observe the bookkeeping against
\Cref{thm:ch11-defect-stokes}.  At the vertex, the outward normal of
$\sigma_j$ is $n_{\sigma_j,\tau}=-u_j$, so
$\mathbf N(\tau)=-\sum_ju_j=0$; the interior term of
\eqref{eq:ch11-tripod-sum} is $-\langle a(0),\sum_ju_j\rangle
=\langle a(0),\mathbf N(\tau)\rangle$, exactly the contraction predicted by
\eqref{eq:ch11-defect-stokes} with $n=1$.  At the endpoint $p_j$ the
outward normal is $+u_j$ and the defect is $u_j$, giving
$\langle a(p_j),u_j\rangle$.  If $\eta$ is supported in a neighbourhood of
the vertex, all three endpoint terms vanish and $\int_C\dprime\eta=0$: the
truncated tripod is balanced at its vertex, which is all that
\Cref{thm:balanced-stokes} requires of an interior cell.
\end{example}

\begin{example}[An unbalanced tripod: Stokes fails]
\label{ex:ch11-tripod-unbalanced}
Keep the notation of the previous example but replace $u_3$ by
$u_3'=(1,1)$, so that the three primitive directions are $(1,0)$, $(0,1)$,
$(1,1)$, all with weight one, and the defect at the vertex is
\[
 \mathbf N(\tau)=-\bigl((1,0)+(0,1)+(1,1)\bigr)=(-2,-2).
\]
Choose $\eta=\chi(x)\,\ddprime x_1$ with $\chi$ a bump function equal to $1$
near the origin and supported well inside the truncation.  Then all endpoint
terms vanish and \eqref{eq:ch11-tripod-sum} gives
\[
 \int_C\dprime\eta=\bigl\langle a(0),\mathbf N(\tau)\bigr\rangle=-2\neq0
\]
with $a=(\chi,0)$.  Stokes' formula therefore fails on this complex, and it
fails by exactly the defect.  A second choice, $\eta=\chi\,\ddprime x_2$,
gives $-2$ as well, and the two computations together determine the whole
defect vector; this is Step 3 of the proof of
\Cref{thm:ch11-balancing-iff-stokes} in a concrete case.

Two remarks.  First, the failure is not repaired by rescaling the weights
\emph{positively}.  The balancing condition
$m_1(1,0)+m_2(0,1)+m_3(1,1)=0$ is the pair of equations $m_1+m_3=0$ and
$m_2+m_3=0$, which leaves $m_3$ free: signed balancings do exist, for
instance $(m_1,m_2,m_3)=(-1,-1,1)$.  What fails is balancing with all three
weights positive, since then the first coordinate $m_1+m_3$ is strictly
positive.  Three rays in an open half-plane therefore admit no balancing
\emph{with positive weights}, which is exactly the assertion of the
vertex-equation example of \Cref{chap:tropical-spaces}; the present
computation verifies it and shows that the positivity hypothesis there
cannot be dropped.
Second, the same configuration is a perfectly good \emph{support}: one can
integrate top-degree forms over it, and the integral is subdivision
invariant by \Cref{prop:integral-subdivision-independent}.  What one cannot do
is integrate by parts.  Balancing is a hypothesis about differentiation, not
about measure.
\end{example}

\begin{example}[Weights are not determined by the support]
\label{ex:ch11-weights-not-support}
Let $C$ be the union of the two coordinate axes in $\mathbf R^2$, a
one-dimensional polyhedral complex with four rays
$\pm e_1,\pm e_2$ from the origin.  Assigning the weight $a$ to both
horizontal rays and $b$ to both vertical rays gives
$\mathbf N(0)=-(a e_1-ae_1+be_2-be_2)=0$ for all $a,b$, so the complex is
balanced for every choice, and each choice gives a different tropical
cycle with the same support.  Consequently
\[
 \int_{(C,(a,b))}\alpha
 =a\int_{\mathbf R\times\{0\}}\alpha+b\int_{\{0\}\times\mathbf R}\alpha
\]
genuinely depends on $(a,b)$.  This is the precise sense in which the
underlying subset of a tropical cycle is insufficient data, as announced at
the start of \Cref{chap:tropical-spaces}.  It also shows that the balancing
condition, being linear in the weights, defines a vector space of cycles on
a given support, of dimension $2$ in this example.
\end{example}

\subsection{From weighted complexes to calibrated tropical spaces}
\label{sec:ch11-global-stokes}

The passage from a single balanced complex to a general calibrated space is
the one place where an analytic input is used.  We state it, indicate how
the cellwise computation of this chapter is deployed, and identify precisely
what is imported.

\begin{theorem}\label{thm:cld-global-stokes}
Let $S$ be a pure $n$-dimensional parlinear space equipped with a calibration
and with empty boundary in the sense of Chambert--Loir--Ducros. For compactly
supported forms $\eta\in A^{n-1,n}_c(S)$ and
$\theta\in A^{n,n-1}_c(S)$,
\[
                              \int_S\dprime\eta=0,
                     \qquad   \int_S\ddprime\theta=0.
\]
With boundary, the corresponding integrals equal the calibrated boundary
integrals.
\end{theorem}

\begin{proof}
This is \cite[Proposition~5.11.1]{CLD25}.  Its proof chooses a locally
finite calibrated chart cover, a subordinate smooth partition of unity, and
compatible polyhedral refinements of the chart representatives.  Compact
support reduces the calculation to finitely many chart terms.  On every
refined chart, \Cref{thm:balanced-stokes} is the cellwise identity: internal
codimension-one contributions cancel by harmony of the calibration, while
the remaining terms are precisely the calibrated boundary integral.  The
existence and compatibility of those chart data are the imported analytic
input; the local sign is the one fixed by
\eqref{eq:superform-boundary-traces}.
\end{proof}

\begin{warning}{Where the analytic input occurs}
The divergence theorem is not the deep step.  The substantive CLD input is
that parlinear charts, calibrations, and their boundaries admit the compatible
refinements used in the proof.  A single preferred tropicalization would not
suffice.
\end{warning}

\section{Green formulae}
\label{sec:ch11-green}

Stokes' formula is a statement about one form; Green's formula is the
statement obtained by applying it to a skew-symmetrized product of two.
Everything in this section is formal once one identity between wedge
products of symmetric forms is available, and we prove that identity first
because its sign is the last one the chapter has to fix.

\begin{lemma}
\label{lem:ch11-mixed-cancellation}
Let $\alpha\in A^{p,p}$ and $\beta\in A^{q,q}$ be symmetric, with
$p+q=n-1$.  Then
\[
 \dprime\alpha\wedge\ddprime\beta
 =-\,\ddprime\alpha\wedge\dprime\beta ,
\]
and consequently, setting
$\omega=\alpha\wedge\ddprime\beta-\ddprime\alpha\wedge\beta\in A^{n-1,n}$,
\[
 \dprime\omega
 =\alpha\wedge\dprimeddprime\beta-\dprimeddprime\alpha\wedge\beta .
\]
\end{lemma}

\begin{proof}
Recall from \Cref{chap:superforms} that $J$ is an algebra automorphism with
$J^2=\operatorname{id}$, that $\ddprime=J\dprime J$, and that a form is
symmetric of bidegree $(p,p)$ when $J\alpha=(-1)^p\alpha$.  Since
$\dprime\alpha\wedge\ddprime\beta$ has bidegree
$(p+1+q,\,p+q+1)=(n,n)$ and $J\tau_n=(-1)^n\tau_n$---each of the $n$
transpositions $\dprime x_i\wedge\ddprime x_i\mapsto
\ddprime x_i\wedge\dprime x_i$ contributes a factor $-1$---we have
\[
 J\bigl(\dprime\alpha\wedge\ddprime\beta\bigr)
 =(-1)^n\,\dprime\alpha\wedge\ddprime\beta .
\]
On the other hand, using that $J$ is multiplicative,
\[
 J\dprime\alpha=\ddprime(J\alpha)=(-1)^p\ddprime\alpha,
 \qquad
 J\ddprime\beta=\dprime(J\beta)=(-1)^q\dprime\beta,
\]
so that
$J(\dprime\alpha\wedge\ddprime\beta)
=(-1)^{p+q}\ddprime\alpha\wedge\dprime\beta
=(-1)^{n-1}\ddprime\alpha\wedge\dprime\beta$.  Comparing the two
evaluations gives the first identity.

For the second, apply the Leibniz rule of
\Cref{prop:superform-differential-identities}.  As $\alpha$ has total degree
$2p$, which is even,
\[
 \dprime(\alpha\wedge\ddprime\beta)
 =\dprime\alpha\wedge\ddprime\beta+\alpha\wedge\dprimeddprime\beta ,
\]
while $\ddprime\alpha$ has total degree $2p+1$, which is odd, so
\[
 \dprime(\ddprime\alpha\wedge\beta)
 =\dprime\ddprime\alpha\wedge\beta-\ddprime\alpha\wedge\dprime\beta
 =\dprimeddprime\alpha\wedge\beta-\ddprime\alpha\wedge\dprime\beta .
\]
Subtracting and using the first identity to cancel
$\dprime\alpha\wedge\ddprime\beta+\ddprime\alpha\wedge\dprime\beta=0$ gives
the assertion.
\end{proof}

\begin{theorem}
\label{thm:ch11-green-defect}
Let $(\mathscr C,m)$ be a pure $n$-dimensional weighted rational polyhedral
complex in $N_{\mathbf R}$, not assumed balanced, and let
$\alpha\in A^{p,p}$, $\beta\in A^{q,q}$ be symmetric of class $C^2$ with
$p+q=n-1$ and with
$\Supp\alpha\cap\Supp\beta\cap|\mathscr C|$ compact.  Then
\begin{equation}
\label{eq:ch11-green-defect}
 \int_{(\mathscr C,m)}
   \bigl(\alpha\wedge\dprimeddprime\beta-\dprimeddprime\alpha\wedge\beta\bigr)
 =(-1)^{n-1}\sum_{\tau\in\mathscr C_{n-1}}\int_\tau
   \Bigl(\iota''_{\mathbf N(\tau)}
     \bigl(\alpha\wedge\ddprime\beta-\ddprime\alpha\wedge\beta\bigr)
   \Bigr)\Big|_\tau .
\end{equation}
If $(\mathscr C,m)$ is balanced the right-hand side vanishes.
\end{theorem}

\begin{proof}
Put $\omega=\alpha\wedge\ddprime\beta-\ddprime\alpha\wedge\beta$, an
$(n-1,n)$-form supported in $\Supp\alpha\cap\Supp\beta$, hence compactly
supported on $|\mathscr C|$.  By \Cref{lem:ch11-mixed-cancellation},
$\dprime\omega$ is the integrand on the left.  Now apply
\Cref{thm:ch11-defect-stokes} to $\omega$.  The last assertion is
\Cref{thm:balanced-stokes}, or directly $\mathbf N(\tau)=0$.
\end{proof}

\begin{corollary}
\label{cor:ch11-green-symmetry}
Let $S$ be a pure $n$-dimensional calibrated parlinear space with empty
boundary, and let $\alpha\in A^{n-1,n-1}(S)$ satisfy
$\dprime\alpha=\ddprime\alpha=0$.  If $f,g$ are smooth and at least one is
compactly supported, then
\[
\int_S f\,\dprimeddprime g\wedge\alpha
=\int_S g\,\dprimeddprime f\wedge\alpha
\]
The first integration by parts below uses $\dprime\alpha=0$ and the second
uses $\ddprime\alpha=0$.
\end{corollary}
\begin{proof}
Stokes applied first to $f\,\ddprime g\wedge\alpha$ gives
\[
 \int_S f\,\dprimeddprime g\wedge\alpha
 =-\int_S \dprime f\wedge \ddprime g\wedge\alpha.
\]
Apply $\ddprime$-Stokes to $g\,\dprime f\wedge\alpha$.  Since
$\ddprime\dprime=-\dprimeddprime$, it gives
\[
 \int_S \dprime f\wedge \ddprime g\wedge\alpha
 =-\int_S g\,\dprimeddprime f\wedge\alpha.
\]
Combining the two identities proves the assertion. If $S$ has boundary, the
same two applications of Stokes retain their boundary terms and give Green's
formula.
\end{proof}

\begin{corollary}
\label{cor:tropical-dirichlet-sign}
Under the hypotheses of \Cref{cor:ch11-green-symmetry}, suppose in addition
that the calibration is positive, that $\alpha$ is strongly positive, and
that $f$ has compact support. Then
\[
 \int_S f\,\dprimeddprime f\wedge\alpha
 =-\int_S\dprime f\wedge\ddprime f\wedge\alpha\leq0.
\]
\end{corollary}
\begin{proof}
The equality is the first integration-by-parts identity with $g=f$.  The
form $\dprime f\wedge\ddprime f$ is strongly positive, and wedging with
$\alpha$ gives a nonnegative top form.
\end{proof}

\begin{example}[The graph energy]
\label{ex:ch11-graph-energy}
On a compact metric graph, take $f$ piecewise smooth and affine near each
vertex.  The right-hand side of
\Cref{cor:tropical-dirichlet-sign} is the negative of the sum of
$\int_e|\partial_s f|^2\,\rd s$ over the edges.  The boundary terms at a
vertex are the sums of outgoing slopes.  They cancel for a harmonic function
and produce the atomic graph Laplacian otherwise.  Thus the polyhedral Green
formula contains the usual graph integration by parts as its
one-dimensional case.
\end{example}

\subsection{The one-dimensional case: metric graphs}
\label{sec:ch11-graph-green}

The last example is the shortest possible statement of a comparison that
deserves to be carried out completely, because \Cref{sec:metric-graphs} fixed
a Laplacian sign convention that the rest of the book uses constantly, and
because that convention is \emph{not} the one a careless reading of the
boundary traces would produce.  This subsection is the sign audit.

Let $\Gamma$ be a compact metric graph in the sense of
\Cref{sec:metric-graphs}, with a model whose edges $e$ are isometrically
identified with intervals $[0,\ell_e]$.  We regard $\Gamma$ as a
one-dimensional polyhedral space with all weights equal to $1$: each edge
carries the lattice for which the arclength coordinate $s$ is a lattice
coordinate, so that $\omega_e=\rd s$ and $\tau_1=\dprime s\wedge\ddprime s$
on $e$.  This is legitimate: the two possible arclength coordinates on $e$
differ by $s\mapsto\ell_e-s$, and $\dprime s\wedge\ddprime s$ is unchanged,
by \Cref{lem:ch11-orientation-drops-out}.  For a function $u$ on $\Gamma$
which is $C^2$ on each edge we write, on the interior of $e$,
\[
 \dprime u=u'\,\dprime s,\qquad
 \ddprime u=u'\,\ddprime s,\qquad
 \dprimeddprime u=u''\,\dprime s\wedge\ddprime s=u''\,\tau_1 .
\]

\begin{lemma}
\label{lem:ch11-graph-trace}
Let $p$ be an endpoint of an edge $e$ and let $\nu\in T_p\Gamma$ be the
tangent direction of $e$ at $p$, in the notation of
\eqref{eq:ch08-outgoing-slope}.  Then, for $u$ of class $C^1$ on $e$ up to
$p$,
\[
 \operatorname{tr}'_{e,p}\bigl(\ddprime u\bigr)
 =-\,d_\nu u(p)\,\tau_0
 \qquad\text{and}\qquad
 \int_{\{p\}}\operatorname{tr}'_{e,p}\bigl(\ddprime u\bigr)=-\,d_\nu u(p).
\]
\end{lemma}

\begin{proof}
Identify $e$ with $[0,\ell_e]$ so that $p$ corresponds to $s=0$; this is
possible after replacing $s$ by $\ell_e-s$, which changes neither
$\tau_1$ nor the assertion.  In the notation of
\Cref{cor:ch11-trace-flux}, $\ddprime u=u'\ddprime s$ lies in
$A^{0,1}=A^{n-1,n}$ with $n=1$, so it corresponds to the vector field
$F_\eta=u'\,\partial_s$ of the first isomorphism there.  The face $\{p\}$ has
$N_{\{p\}}=0$, and the outward primitive normal of $e$ at $p$ is
$n_{e,p}=-\partial_s$, since $\partial_s$ points into $e$.  Hence
$F_\eta=u'(0)\partial_s=(-u'(0))\,n_{e,p}$
and $\lambda=-u'(0)$.  Finally $d_\nu u(p)=u'(0)$ by
\eqref{eq:ch08-outgoing-slope}, because the unit-speed geodesic
representing $\nu$ is $s\mapsto s$.  The lattice measure of a point is the
unit mass, whence the second formula.
\end{proof}

\begin{warning}{Outgoing slopes point inward}
The outgoing slope $d_\nu u(p)$ of \Cref{sec:metric-graphs} is the
derivative of $u$ in the direction leading \emph{away} from $p$ along the
edge, that is, in the direction of the \emph{inward} normal $u_{e/p}$ of the
edge along its face $\{p\}$.  The boundary trace of the present chapter uses
the \emph{outward} normal $n_{e,p}=-u_{e/p}$.  The two conventions therefore
differ by a sign, and \Cref{lem:ch11-graph-trace} is that sign.  A reader who
defines the graph Laplacian as the sum of outward boundary traces obtains
$-\Delta_\Gamma$, and every subsequent statement---positivity of a curvature
measure at a convex corner, the sign of the Dirichlet form, the direction of
the maximum principle---is reversed.  This is the promised place where the
naive orientation convention gives the wrong sign, and it is worth stressing
that no dimension-two test detects it: it is a one-dimensional error, of the
kind that \Cref{ex:ch11-square-wrong-sign} does not see, just as
\Cref{ex:ch11-square-wrong-sign} is invisible in dimension one.
\end{warning}

\begin{theorem}
\label{thm:ch11-graph-green}
Let $\Gamma$ be a compact metric graph and let $f,g\in\PL(\Gamma)$.  Fix a
model of $\Gamma$ in which both are affine on every edge.  Then:
\begin{enumerate}
\item\label{itm:ch11-graph:lap}
  For every $p\in\Gamma$,
  \[
   \bigl(\Delta_\Gamma f\bigr)(\{p\})
   =-\sum_{e\ni p}\int_{\{p\}}\operatorname{tr}'_{e,p}\bigl(\ddprime f\bigr)
   =-\int_{\{p\}}\operatorname{tr}'_{\Gamma,p}\bigl(\ddprime f\bigr),
  \]
  where the last expression abbreviates the sum of the traces of all edges
  at $p$.  Equivalently, $\Delta_\Gamma f$ is the negative of the total
  outward flux of $\ddprime f$.
\item\label{itm:ch11-graph:mass}
  $\Delta_\Gamma f(\Gamma)=0$; this is \eqref{eq:lap-mass-zero}.
\item\label{itm:ch11-graph:parts}
  $\displaystyle\int_\Gamma f\,\rd\Delta_\Gamma g
   =-\int_\Gamma\dprime f\wedge\ddprime g$, where the right-hand side is the
  sum of the edge integrals of $f'g'\,\rd s$; this is the second equality of
  \eqref{eq:graph-green}.
\item\label{itm:ch11-graph:sym}
  $\displaystyle\int_\Gamma f\,\rd\Delta_\Gamma g
   =\int_\Gamma g\,\rd\Delta_\Gamma f$; this is the first equality of
  \eqref{eq:graph-green}.
\end{enumerate}
\end{theorem}

\begin{proof}
Throughout we use the following edgewise form of Stokes.  Each edge $e$ is a
one-dimensional rational polytope in $\mathbf R$ with $N_e=\mathbf Z$, so
\Cref{thm:polyhedral-stokes} applies to it and reads
\begin{equation}
\label{eq:ch11-graph-edge-stokes}
 \int_e\dprime\eta
 =\sum_{p\in\partial e}\int_{\{p\}}\operatorname{tr}'_{e,p}(\eta),
 \qquad \eta\in A^{0,1}(e).
\end{equation}
Summing \eqref{eq:ch11-graph-edge-stokes} over the edges of the model and
grouping the right-hand side by vertices is the only global step used below;
no ambient affine space is required, because both sides are computed inside
individual edges.

\cref{itm:ch11-graph:lap} is \Cref{lem:ch11-graph-trace} summed over the
edges at $p$, together with the definition
\eqref{eq:graph-laplacian} of $\Delta_\Gamma$ as the sum of outgoing slopes.
For a point $p$ interior to an edge the two traces are $-f'$ and $f'$ and
cancel, in accordance with
\Cref{lem:ch08-laplacian-well-defined}\,(1).

\cref{itm:ch11-graph:mass}.  Take $\eta=\ddprime f$ in
\eqref{eq:ch11-graph-edge-stokes}.  Since $f$ is affine on each edge,
$\dprime\eta=f''\tau_1=0$ there, so every left-hand side vanishes.  Summing
over edges and grouping by vertices, the total right-hand side is
$\sum_p\sum_{e\ni p}\int_{\{p\}}\operatorname{tr}'_{e,p}(\ddprime f)$,
which by \cref{itm:ch11-graph:lap} equals $-\Delta_\Gamma f(\Gamma)$.

\cref{itm:ch11-graph:parts}.  Take $\eta=f\,\ddprime g$ in
\eqref{eq:ch11-graph-edge-stokes}.  By the Leibniz rule
$\dprime\eta=\dprime f\wedge\ddprime g+f\,\dprimeddprime g$, and the second
term vanishes edgewise because $g$ is affine.  On each edge
$\dprime f\wedge\ddprime g=f'g'\,\tau_1$, so the sum of the left-hand sides
is $\int_\Gamma f'g'\,\rd s$.  The sum of the right-hand sides is
$\sum_pf(p)\sum_{e\ni p}\int_{\{p\}}\operatorname{tr}'_{e,p}(\ddprime g)$,
where $f$ could be taken outside because it is continuous and the trace is
evaluated at the single point $p$.  By \cref{itm:ch11-graph:lap} this is
$-\sum_pf(p)\Delta_\Gamma g(\{p\})=-\int_\Gamma f\,\rd\Delta_\Gamma g$.
Rearranging gives the assertion.

\cref{itm:ch11-graph:sym} follows from \cref{itm:ch11-graph:parts} and the
symmetry of $(f,g)\mapsto\int_\Gamma f'g'\,\rd s$.  Alternatively it is the
edgewise Green formula, obtained by taking
$\eta=f\,\ddprime g-g\,\ddprime f$ in \eqref{eq:ch11-graph-edge-stokes}: the
left-hand sides vanish because $f$ and $g$ are affine on each edge, and the
right-hand side is
$\sum_p\bigl(g(p)\Delta_\Gamma f(\{p\})-f(p)\Delta_\Gamma g(\{p\})\bigr)$ up
to the global sign of \cref{itm:ch11-graph:lap}.
\end{proof}

Thus \Cref{prop:graph-green} is exactly the one-dimensional case of the
polyhedral Green formula, once the sign of
\Cref{lem:ch11-graph-trace} is inserted.  We record the two sign checks that
were the point of the exercise.

\begin{example}[A segment, checked against the metric-graph conventions]
\label{ex:ch11-graph-segment-check}
Let $\Gamma=[0,\ell]$ and $f(t)=t$.  According to
\eqref{eq:graph-laplacian}, $\Delta_\Gamma f=\delta_0-\delta_\ell$: at $0$
the unique outgoing slope is $+1$, at $\ell$ it is $-1$.  According to
\Cref{lem:ch11-graph-trace}, the outward traces are
$\operatorname{tr}'_{\Gamma,0}(\ddprime f)=-1$ and
$\operatorname{tr}'_{\Gamma,\ell}(\ddprime f)=+1$, so their negatives are
$+1$ and $-1$, in agreement.  As a cross-check,
\Cref{ex:ch11-segment-stokes} with $m=1$ and $\eta=\ddprime f=\ddprime t$
gives $\int_\Gamma\dprime\ddprime f=f'(\ell)-f'(0)=0$ and total defect mass
$-\Delta_\Gamma f(\Gamma)=0$, consistent with
\eqref{eq:lap-mass-zero}.  Finally
$\int_\Gamma f\,\rd\Delta_\Gamma f=0-\ell=-\ell$ and
$-\int_0^\ell|f'|^2\,\rd t=-\ell$; this is the computation recorded after
\Cref{prop:graph-green}, and it also confirms
\Cref{cor:tropical-dirichlet-sign} in the case $n=1$, $\alpha=1$: the
Dirichlet pairing is negative semidefinite, as in
\Cref{cor:ch08-dirichlet-energy}.
\end{example}

\begin{example}[A tripod, checked against the model computation of Chapter~8]
\label{ex:ch11-graph-tripod-check}
Let $\Gamma$ be the tripod with central vertex $c$ and three legs of length
one, with tips $q_0,q_1,q_2$.  Fix $a_0,a_1\in\mathbf R$ and let $\varphi$
be the piecewise affine function vanishing at $c$ and on the third leg,
affine on each leg, with $\varphi(q_i)=a_i$ for $i=0,1$.  Then the outgoing
slopes at $c$ are $a_0,a_1,0$, so
$\Delta_\Gamma\varphi=(a_0+a_1)\delta_c-a_0\delta_{q_0}-a_1\delta_{q_1}$,
which is the graph computation appearing in \Cref{ex:ch08-p1-tripod}.

Now compute the same measure by boundary traces.  On the $i$th leg,
parametrized by arclength $s\in[0,1]$ from $c$, one has
$\ddprime\varphi=a_i\,\ddprime s$ for $i=0,1$ and $0$ on the third leg.  At
$c$ the outward normal of each leg is $-\partial_s$, so the trace of the
$i$th leg is $-a_i$ and the total outward flux at $c$ is $-(a_0+a_1)$; its
negative is $a_0+a_1$, matching.  At $q_i$ the outward normal is
$+\partial_s$, the trace is $a_i$, and its negative is $-a_i$, again
matching.  The total mass is $0$.

The tropical content is visible: at $c$ the three primitive inward
directions, viewed in a chart identifying a neighbourhood of $c$ with the
tropical line of \Cref{fig:tropical-line-balancing}, sum to zero, so
$\mathbf N(c)=0$ and $\Gamma$ is balanced at $c$; whereas at each tip
$\mathbf N(q_i)$ is a primitive vector and $\Gamma$ has genuine boundary
there.  Correspondingly $\Delta_\Gamma\varphi$ may be nonzero at $c$ only
because $\varphi$ is not smooth in the chart, while its atoms at the tips
are boundary terms.  In \Cref{ex:ch08-p1-tripod} the same three atoms
appeared as $\Delta_\Gamma\varphi_D$, the residual degrees of the model line
bundle being instead the atoms of
$\delta_{x_{\widetilde C}}+\Delta_\Gamma\varphi_D$, which differ in the
central atom by the mass $1$ of $\delta_{x_{\widetilde C}}$; the two
Laplacian computations agree atom by atom, and the mass ledger of
\Cref{sec:ch08-mass-ledger} records the total.
\end{example}

\begin{example}[The Green function of a segment]
\label{ex:ch11-green-function-check}
For $q\in(0,1)$ let $g_q$ be the function of
\eqref{eq:segment-green} on $\Gamma=[0,1]$, so that
$\Delta g_q=(1-q)\delta_0+q\delta_1-\delta_q$.  By
\Cref{thm:ch11-graph-green}\,\cref{itm:ch11-graph:parts},
\[
 \int_\Gamma g_q\,\rd\Delta g_q=-\int_0^1|g_q'|^2\,\rd s=-q(1-q),
\]
and the left-hand side computed from the atoms is
$(1-q)g_q(0)+qg_q(1)-g_q(q)=0+0-q(1-q)$.  The two agree.  This is the sign
test announced in \Cref{cor:tropical-dirichlet-sign}: the quantity
$\int f\,\dprimeddprime f$ is $\le0$, while the Dirichlet energy
$\int|f'|^2$ is $\ge0$.  Reversing the trace convention would produce a
Green function of the wrong sign and a negative Dirichlet energy.
\end{example}

\section{Pushforward and the projection formula}
\label{sec:ch11-pushforward}

For an integer-affine map $F:N_{\mathbf R}\to N'_{\mathbf R}$ that is
injective on a top cell, the pushforward weight is multiplied by the lattice
index $[N'_{F(\sigma)}:F(N_\sigma)]$. This is forced by change of variables.
We now prove the change-of-variables and projection formulas in the
generality needed later, allowing a facet to be collapsed by $F$.  Properness
is the condition that makes the image complex locally finite and every
compactly supported integral finite.

\subsection{Pushforward with collapsing cells}
\label{sec:ch11-pushforward-def}

\begin{definition}
\label{def:ch11-proper-pushforward}
\index{pushforward!of a tropical cycle}
Let $(\mathscr C,m)$ be a pure $n$-dimensional weighted rational polyhedral
complex in $N_{\mathbf R}$ and let $F:N_{\mathbf R}\to N'_{\mathbf R}$ be
integer affine, that is, affine with linear part carrying $N$ into $N'$.
Assume $F$ is \emph{proper on $|\mathscr C|$}: the preimage in
$|\mathscr C|$ of every compact subset of $N'_{\mathbf R}$ is compact.  Call
a facet $\sigma$ \emph{nondegenerate} if $\dim F(\sigma)=n$, equivalently if
$F$ is injective on $N_{\sigma,\mathbf R}$, and \emph{collapsed} otherwise.
After a subdivision for which the images of the nondegenerate facets form a
polyhedral complex $\mathscr C'$ in $N'_{\mathbf R}$, set, for
$\sigma'\in\mathscr C'_n$,
\[
 m'_{\sigma'}
 =\sum_{\substack{\sigma\in\mathscr C_n\ \text{nondegenerate}\\
                  F(\sigma)=\sigma'}}
   m_\sigma\,[N'_{\sigma'}:F(N_\sigma)],
\]
and let $F_*(\mathscr C,m)$ be the weighted complex obtained by discarding
the cells of weight zero.  Collapsed facets contribute nothing.
\end{definition}

The subdivision exists and the resulting complex is locally finite precisely
because $F$ is proper: a compact subset of $N'_{\mathbf R}$ meets only
finitely many image cells, since it pulls back to a compact subset of
$|\mathscr C|$, which meets only finitely many cells.  When every facet is
nondegenerate, the weight formula is the one of
\Cref{chap:tropical-spaces}, and
\Cref{prop:tropical-pushforward-balanced} applies.

\begin{lemma}
\label{lem:ch11-cell-change-of-variables}
Let $\sigma$ be an $n$-dimensional rational polyhedron in $N_{\mathbf R}$,
let $F$ be integer affine, and let $\alpha$ be an $(n,n)$-form on
$N'_{\mathbf R}$ whose support meets $F(\sigma)$ in a compact set.
\begin{enumerate}
\item\label{itm:ch11-cov:collapse}
  If $\sigma$ is collapsed, then $(F^*\alpha)|_\sigma=0$.
\item\label{itm:ch11-cov:nondeg}
  If $\sigma$ is nondegenerate, with $\sigma'=F(\sigma)$ and
  $d=[N'_{\sigma'}:F(N_\sigma)]$, then
  \[
   \int_\sigma F^*\alpha=d\int_{\sigma'}\alpha .
  \]
\end{enumerate}
\end{lemma}

\begin{proof}
\cref{itm:ch11-cov:collapse}
Let $B$ be the linear part of $F$.  If $\sigma$ is collapsed then
$B|_{N_{\sigma,\mathbf R}}$ has rank $<n$, so its image has dimension
$k<n$.  A pullback $F^*\alpha$ evaluated on $n$ primed and $n$
double-primed vectors tangent to $\sigma$ is $\alpha$ evaluated on their
images, which span a space of dimension $k<n$; an alternating form of degree
$n$ in each block vanishes on such a family.  Hence $(F^*\alpha)|_\sigma=0$.

\cref{itm:ch11-cov:nondeg}
Choose lattice bases of $N_\sigma$ and of $N'_{\sigma'}$ and let
$D$ be the matrix of $B|_{N_{\sigma,\mathbf R}}$ in these bases, so that
$|\det D|=d$.  Writing $\alpha=g\,\tau'_n$ in the target coordinates, the
computation of \Cref{lem:ch11-orientation-drops-out} gives
$F^*\alpha=(g\circ F)\det(D)^2\,\tau_n$ on the span of $\sigma$, whence
\[
 \int_\sigma F^*\alpha
 =d^2\int_\sigma (g\circ F)\,\omega_\sigma
 =d^2\cdot d^{-1}\int_{\sigma'}g\,\omega_{\sigma'}
 =d\int_{\sigma'}\alpha ,
\]
the middle equality being the affine change-of-variables formula together
with $F_*\omega_\sigma=d^{-1}\omega_{\sigma'}$, which holds because the
image of a lattice basis of $N_\sigma$ spans a parallelotope of
$\omega_{\sigma'}$-volume $d$.
\end{proof}

\begin{warning}{Two factors of the index, one removed by the measure}
The pullback of a top-degree superform acquires $\det(D)^2$, because both
labelled blocks are pulled back; the change of variables for the measure
removes one factor $|\det D|$.  What remains is one factor $d$, and that is
the pushforward weight.  This is the content of
\Cref{ex:ch11-nonunimodular}, and it is the reason that the tropical
multiplicity is a degree and not its square.  For an ordinary $n$-form only
one factor would appear and it would be cancelled entirely, which is why
ordinary forms cannot see tropical multiplicities.
\end{warning}

\begin{proposition}
\label{prop:ch11-pushforward-integral}
Let $(\mathscr C,m)$ and $F$ be as in
\Cref{def:ch11-proper-pushforward}, and let $\alpha$ be an $(n,n)$-form with
compact support on $N'_{\mathbf R}$.  Then $F^*\alpha$ has compact support
on $|\mathscr C|$ and
\[
 \int_{(\mathscr C,m)}F^*\alpha
 =\int_{F_*(\mathscr C,m)}\alpha .
\]
\end{proposition}

\begin{proof}
Compactness of $\Supp(F^*\alpha)\cap|\mathscr C|$ is properness.  Refine so
that \Cref{def:ch11-proper-pushforward} applies; by
\Cref{prop:integral-subdivision-independent} neither side changes under
refinement.  By \Cref{lem:ch11-cell-change-of-variables}, collapsed facets
contribute zero and
\[
 \int_{(\mathscr C,m)}F^*\alpha
 =\sum_{\sigma\ \text{nondeg.}}m_\sigma\int_\sigma F^*\alpha
 =\sum_{\sigma\ \text{nondeg.}}
   m_\sigma[N'_{F(\sigma)}:F(N_\sigma)]\int_{F(\sigma)}\alpha .
\]
Grouping the terms according to the image cell $\sigma'=F(\sigma)$ produces
exactly $\sum_{\sigma'}m'_{\sigma'}\int_{\sigma'}\alpha$, which is the
right-hand side.  All sums are finite by properness.
\end{proof}

\begin{theorem}
\label{thm:ch11-pushforward-balanced}
Let $(\mathscr C,m)$ be balanced and let $F$ be integer affine and proper on
$|\mathscr C|$.  Then $F_*(\mathscr C,m)$ is balanced.
\end{theorem}

\begin{proof}
Write $(\mathscr C',m')=F_*(\mathscr C,m)$ and let $\eta$ be a smooth
$(n-1,n)$-form with compact support on $N'_{\mathbf R}$; by the last
assertion of \Cref{thm:ch11-balancing-iff-stokes} it suffices to test
against such forms.  Since the definition of $F_*$ used
only the images of cells, $|\mathscr C'|\subset F(|\mathscr C|)$, so
$F^*\eta$ is a smooth $(n-1,n)$-form whose support meets $|\mathscr C|$ in a
compact set, by properness.  Pullback along an affine map commutes with
$\dprime$ (\Cref{chap:superforms}), so
$\dprime(F^*\eta)=F^*(\dprime\eta)$.  Therefore, by
\Cref{prop:ch11-pushforward-integral} applied to $\alpha=\dprime\eta$ and by
\Cref{thm:balanced-stokes} applied to $F^*\eta$ on $(\mathscr C,m)$,
\[
 \int_{(\mathscr C',m')}\dprime\eta
 =\int_{(\mathscr C,m)}F^*(\dprime\eta)
 =\int_{(\mathscr C,m)}\dprime(F^*\eta)
 =0 .
\]
As $\eta$ was arbitrary, criterion
\cref{itm:ch11-bal:dprime} of
\Cref{thm:ch11-balancing-iff-stokes} shows that
$(\mathscr C',m')$ is balanced.
\end{proof}

\begin{remark}[A proof by analysis of a combinatorial statement]
\label{rem:ch11-pushforward-proof-style}
\Cref{thm:ch11-pushforward-balanced} is a statement about lattice indices and
primitive normal vectors, and \Cref{chap:tropical-spaces} proved it directly
in the case where no facet is collapsed
(\Cref{prop:tropical-pushforward-balanced}).  The proof just given instead
deduces it from the analytic characterization of balancing.  This is not
merely an economy: the combinatorial proof in the collapsing case requires a
separate discussion of the cells of $\mathscr C$ of dimension $n-1$ that
dominate a codimension-one cell of the image, whereas the analytic proof is
insensitive to that distinction because $\dprime$ and pullback commute
regardless.  The same technique will reappear whenever a tropical operation
has to be shown to preserve balancing.  A purely combinatorial treatment,
valid for arbitrary integer-affine maps, is
\cite[\S4.6]{MR18Tropical}.
\end{remark}

\begin{corollary}
\label{cor:ch11-projection-formula}
Let $(\mathscr C,m)$ be a pure $n$-dimensional weighted complex and let $F$
be integer affine and proper on $|\mathscr C|$.  For every smooth
$(p,p)$-form $\beta$ on $N'_{\mathbf R}$ and every $(n-p,n-p)$-form
$\gamma$ with compact support on $N'_{\mathbf R}$,
\[
 \int_{(\mathscr C,m)}F^*\beta\wedge F^*\gamma
 =\int_{F_*(\mathscr C,m)}\beta\wedge\gamma .
\]
Equivalently, in the notation of \Cref{chap:cld-currents} where
$\delta_{(\mathscr C,m)}$ is the integration current,
\[
 F_*\delta_{(\mathscr C,m)}=\delta_{F_*(\mathscr C,m)}
 \qquad\text{and}\qquad
 F_*\bigl(F^*\beta\wedge\delta_{(\mathscr C,m)}\bigr)
 =\beta\wedge F_*\delta_{(\mathscr C,m)} .
\]
\end{corollary}

\begin{proof}
The first identity is \Cref{prop:ch11-pushforward-integral} applied to
$\alpha=\beta\wedge\gamma$, using $F^*(\beta\wedge\gamma)
=F^*\beta\wedge F^*\gamma$.  The identity
$F_*\delta_{(\mathscr C,m)}=\delta_{F_*(\mathscr C,m)}$ is the case $p=0$,
$\beta=1$, read as an equality of functionals on compactly supported test
forms, since $\langle F_*T,\gamma\rangle=\langle T,F^*\gamma\rangle$ by
\Cref{prop:current-pushforward-formal}.  The last identity then follows by
evaluating both sides on a test form $\gamma$: the left-hand side gives
$\int_{(\mathscr C,m)}F^*\beta\wedge F^*\gamma$ and the right-hand side
gives $\int_{F_*(\mathscr C,m)}\beta\wedge\gamma$.
\end{proof}

\subsection{A projection computed by hand}
\label{sec:ch11-projection-example}

\begin{example}[A balanced curve in $\mathbf R^2$ projected to $\mathbf R$]
\label{ex:ch11-projection-curve}
Let $N=\mathbf Z^2$ and let $C$ be the one-dimensional fan with three rays
\[
 \rho_1=\mathbf R_{\ge0}(2,1),\quad
 \rho_2=\mathbf R_{\ge0}(2,-1),\quad
 \rho_3=\mathbf R_{\ge0}(-1,0),
\]
of weights $m_1=m_2=1$ and $m_3=4$.  All three direction vectors are
primitive, and
\[
 1\cdot(2,1)+1\cdot(2,-1)+4\cdot(-1,0)=(0,0),
\]
so $C$ is balanced at the origin, its only codimension-one cell.  Let
$\pi(x,y)=x$, with $N'=\mathbf Z$.

\emph{Properness.}  Parametrize $\rho_1$ by $t\mapsto t(2,1)$ with $t\ge0$;
then $\pi=2t$, so $\pi^{-1}([-R,R])\cap\rho_1$ is the compact set
$0\le t\le R/2$.  Similarly for $\rho_2$, while on $\rho_3$ one has
$\pi=-t$ and the preimage is $0\le t\le R$.  Hence $\pi$ is proper on $|C|$.

\emph{The pushforward.}  No ray is collapsed.  For $\rho_1$,
$N_{\rho_1}=\mathbf Z(2,1)$ and $\pi(N_{\rho_1})=2\mathbf Z$, so the index is
$[\mathbf Z:2\mathbf Z]=2$ and the contribution to $\mathbf R_{\ge0}$ is
$1\cdot2=2$; the same for $\rho_2$.  For $\rho_3$,
$N_{\rho_3}=\mathbf Z(1,0)$ and $\pi(N_{\rho_3})=\mathbf Z$, index $1$, and
the contribution to $\mathbf R_{\le0}$ is $4\cdot1=4$.  Therefore
\[
 \pi_*C=4\,[\mathbf R_{\ge0}]+4\,[\mathbf R_{\le0}]=4\,[\mathbf R],
\]
which is balanced, as \Cref{thm:ch11-pushforward-balanced} predicts:
$4\cdot(+1)+4\cdot(-1)=0$.

\emph{The projection formula, verified.}  Let
$\alpha=g(x)\,\dprime x\wedge\ddprime x$ with $g$ smooth and compactly
supported on $\mathbf R$.  On $\rho_1$, with lattice coordinate $t$, one has
$\pi\circ(t\mapsto t(2,1))=2t$, so
 \[
  \begin{aligned}
  \pi^*\alpha\big|_{\rho_1}
    &=g(2t)\,\dprime(2t)\wedge\ddprime(2t)
      =4\,g(2t)\,\dprime t\wedge\ddprime t,\\
  \int_{\rho_1}\pi^*\alpha
    &=4\int_0^\infty g(2t)\,\rd t
      =2\int_0^\infty g(x)\,\rd x .
  \end{aligned}
 \]
The same computation on $\rho_2$ gives $2\int_0^\infty g$.  On $\rho_3$ the
lattice coordinate $t$ satisfies $x=-t$, so
$\pi^*\alpha|_{\rho_3}=g(-t)\,\dprime t\wedge\ddprime t$ and
$\int_{\rho_3}\pi^*\alpha=\int_{-\infty}^0g(x)\,\rd x$.  Weighting and
summing,
\[
 \int_{(C,m)}\pi^*\alpha
 =1\cdot2\!\int_0^\infty\! g+1\cdot2\!\int_0^\infty\! g
   +4\!\int_{-\infty}^0\! g
 =4\int_{\mathbf R}g
 =\int_{4[\mathbf R]}\alpha ,
\]
in agreement with \Cref{prop:ch11-pushforward-integral}.  Note where the
factor $4$ in the pullback on $\rho_1$ went: two units of it came from the
two labelled copies of $\dprime(2t)$ and $\ddprime(2t)$; the change of
variables $x=2t$ returned one factor $2$; and the surviving factor $2$ is
the lattice index.  This is the arithmetic of the warning above, made
concrete.

\emph{A consistency check by Stokes.}  Take
$\eta=h(x)\,\ddprime x$ with $h$ smooth and compactly supported.  Then
$\pi^*\eta=h(x)\,\ddprime x$ regarded on $\mathbf R^2$, and
$\int_{(C,m)}\dprime\pi^*\eta$ must vanish by
\Cref{thm:balanced-stokes}.  Directly: on $\rho_j$ with primitive direction
$u_j=(a_j,b_j)$ one has $\pi^*\eta|_{\rho_j}=h(a_jt)a_j\,\ddprime t$, so
\[
 \int_{\rho_j}\dprime\pi^*\eta
 =\bigl[h(a_jt)a_j\bigr]_{t=0}^{\infty}=-a_jh(0),
\]
and the weighted sum is $-h(0)\sum_jm_ja_j=-h(0)\cdot0=0$, the vanishing
being exactly the first coordinate of the balancing relation.  Downstairs,
$\int_{4[\mathbf R]}\dprime\eta=4\bigl[h\bigr]_{-\infty}^{\infty}=0$, as it
must be.
\end{example}

\begin{example}[Properness is not decorative]
\label{ex:ch11-pushforward-nonproper}
Let $N=\mathbf Z^2$ and let $C=\bigcup_{k\in\mathbf Z}\mathbf R\times\{k\}$
be the disjoint union of the integral horizontal lines, each a single
one-dimensional cell of weight $1$.  There are no codimension-one cells at
all, so $C$ is balanced vacuously, and it is a locally finite complex.  Let
$\pi(x,y)=x$.  Every line is nondegenerate and maps isomorphically onto
$\mathbf R$ with index $1$, so the naive pushforward weight on $\mathbf R$
would be $\sum_{k\in\mathbf Z}1$, which is not a well-defined integer.
Correspondingly, for $\alpha=g\,\dprime x\wedge\ddprime x$ with
$g\ge0$ not identically zero,
\[
 \int_C\pi^*\alpha=\sum_{k\in\mathbf Z}\int_{\mathbf R}g(x)\,\rd x=+\infty .
\]
The map $\pi$ is not proper on $|C|$: the preimage of $[-1,1]$ is an
infinite disjoint union of segments.  Both defects are genuine, and both are
removed by the properness hypothesis of
\Cref{def:ch11-proper-pushforward}.

By contrast, the tropical line $L$ of
\Cref{fig:tropical-line-balancing} with $\pi(x,y)=x$ is a case where
properness fails---the vertical ray has non-compact fibre over the
origin---and yet the cellwise formula is finite for every compactly supported
top form, because its pullback to the collapsed cell vanishes by
\Cref{lem:ch11-cell-change-of-variables}\,\cref{itm:ch11-cov:collapse}.  One
obtains the formal image cycle $[\mathbf R]$ and
$\int_L\pi^*\alpha=\int_{\mathbf R}g$.  Properness is thus sufficient, not
necessary for this particular calculation; without it, however,
\Cref{def:ch11-proper-pushforward} does not define $\pi_*L$, and there is no
general pushforward statement, as the first half of the example shows.
\end{example}

\begin{remark}[Functoriality]
\label{rem:ch11-pushforward-functorial}
If $F:N_{\mathbf R}\to N'_{\mathbf R}$ and $G:N'_{\mathbf R}\to
N''_{\mathbf R}$ are integer affine, $F$ is proper on $|\mathscr C|$, and
the restriction of $G$ to $F(|\mathscr C|)$ is proper, then $G\circ F$ is
proper on $|\mathscr C|$ and $(G\circ F)_*=G_*\circ F_*$.  In particular,
$G$ is proper on the smaller set $|F_*\mathscr C|$.  For the weights this is the
multiplicativity of lattice indices used in
\Cref{prop:tropical-pushforward-balanced}, together with the observation
that a facet collapsed by $G\circ F$ is either collapsed by $F$ or has
image collapsed by $G$.  Alternatively it follows from
\Cref{prop:ch11-pushforward-integral} together with the observation that a
weighted complex is determined by its integration functional: if
$\int_{(\mathscr C,m)}\alpha=\int_{(\mathscr C,\tilde m)}\alpha$ for all
compactly supported $\alpha$, then choosing $\alpha=\chi\tau_n$ with $\chi$
a bump at an interior point of a facet $\sigma$, and noting that only
$\sigma$ contributes, gives $m_\sigma=\tilde m_\sigma$.
\end{remark}

\section{Tropical intersection products as current products}
\label{sec:ch11-tropical-intersection}

\Cref{chap:tropical-spaces} constructed a stable intersection product of
tropical cycles by a purely combinatorial fan-displacement argument.  The
present chapter has constructed an integration functional attached to a
cycle.  The two constructions must be compatible, and the purpose of this
section is to say in what sense and to prove it in the cases used later.

The difficulty is that the product of two currents has no meaning in
general, so a literal statement of the form
$\delta_X\wedge\delta_Y=\delta_{X\cdot Y}$ is not available inside the smooth
calculus.  What is available, and what we prove, is that the intersection
current is the limit of products in which one factor has been replaced by a
smooth approximation.  This is the exact analogue of the classical fact that
the current of integration over a submanifold is a weak limit of smooth
bump forms concentrated near it, and it suffices for every application in
this book.  The enlarged calculus in which the product of two polyhedral
currents can be formed directly is the delta-form calculus of
\Cref{chap:delta-forms}; see \Cref{thm:delta-preform-algebra}.

\subsection{Integration currents}
\label{sec:ch11-integration-current}

\begin{definition}
\label{def:ch11-integration-current}
\index{current!of integration over a tropical cycle}
Let $(\mathscr C,m)$ be a pure $n$-dimensional weighted rational polyhedral
complex in $N_{\mathbf R}$, $\dim N_{\mathbf R}=r$.  Its \emph{integration
current} is the linear functional
\[
 \bigl\langle\delta_{(\mathscr C,m)},\alpha\bigr\rangle
 =\int_{(\mathscr C,m)}\alpha ,
 \qquad \alpha\in A^{n,n}_c(N_{\mathbf R}),
\]
a current of bidimension $(n,n)$, that is, of bidegree $(r-n,r-n)$, in the
sense of \Cref{chap:cld-currents}.
\end{definition}

By \Cref{prop:integral-subdivision-independent} the functional depends only on
the tropical cycle, not on the chosen complex; and it is continuous for the
test-form topology of \Cref{chap:cld-currents}, since it is bounded by a
constant times the supremum of the coefficients of $\alpha$ on any fixed
compact set.

\begin{theorem}
\label{thm:ch11-current-closed}
Let $(\mathscr C,m)$ be a pure $n$-dimensional weighted rational polyhedral
complex.  Then $(\mathscr C,m)$ is balanced if and only if
$\dprime\delta_{(\mathscr C,m)}=0$, if and only if
$\ddprime\delta_{(\mathscr C,m)}=0$.
\end{theorem}

\begin{proof}
By the sign convention of \Cref{chap:cld-currents}, $\dprime T=0$ means
$\langle T,\dprime\eta\rangle=0$ for every compactly supported test form
$\eta$ of the appropriate bidegree, the prefactor $(-1)^{p+q+1}$ being
invertible.  For $T=\delta_{(\mathscr C,m)}$ this reads
$\int_{(\mathscr C,m)}\dprime\eta=0$ for all
$\eta\in A^{n-1,n}_c$, which is condition
\cref{itm:ch11-bal:dprime} of \Cref{thm:ch11-balancing-iff-stokes}.  The
double-primed statement is condition \cref{itm:ch11-bal:ddprime}.
\end{proof}

This is the promised reformulation.  A tropical cycle is exactly a
polyhedral support with weights whose integration current is closed;
balancing is closedness.  Combined with \Cref{cor:ch11-projection-formula},
it yields the current-level projection formula in the form used in
\Cref{chap:cld-currents,chap:cld-metrized-bundles}.

\subsection{Smoothing the current of an affine cycle}
\label{sec:ch11-smoothing}

\begin{lemma}
\label{lem:ch11-linear-regularization}
Let $A\subset N_{\mathbf R}$ be a rational affine subspace of dimension $k$
with saturated direction lattice $N_A\subset N$, and let
$w\in\mathbf Z\setminus\{0\}$ be a weight.  Choose
$\ell_{k+1},\ldots,\ell_r\in M$ vanishing on $N_A$ and inducing a basis of
$(N/N_A)^*$, and constants $a_{k+1},\ldots,a_r\in\mathbf R$ such that the
integer-affine functions
\[
 \xi_j(x)=\langle\ell_j,x\rangle-a_j
 \qquad (k<j\le r)
\]
satisfy $A=\{\xi_{k+1}=\cdots=\xi_r=0\}$.  Let $\rho\ge0$ be smooth with
compact support on $\mathbf R^{r-k}$ and $\int\rho=1$, put
$\rho_\varepsilon(u)=\varepsilon^{-(r-k)}\rho(u/\varepsilon)$, and set
\[
 \alpha^A_\varepsilon
 =w\,\rho_\varepsilon(\xi_{k+1},\ldots,\xi_r)\,
   \bigwedge_{j=k+1}^r\dprime\xi_j\wedge\ddprime\xi_j .
\]
Then $\alpha^A_\varepsilon$ is a smooth $(r-k,r-k)$-form which is
$\dprime$- and $\ddprime$-closed, is strongly positive when $w>0$, has
support in the set $\{|\xi''|\le C\varepsilon\}$ for some constant $C$, and
satisfies
\[
 \lim_{\varepsilon\to0}\bigl[\alpha^A_\varepsilon\bigr]
 =\delta_{(A,w)}
\]
weakly, that is,
$\int_{N_{\mathbf R}}\alpha^A_\varepsilon\wedge\beta\to\int_{(A,w)}\beta$
for every $\beta\in A^{k,k}_c(N_{\mathbf R})$.
\end{lemma}

\begin{proof}
Complete the linear parts $\ell_{k+1},\ldots,\ell_r$ to a basis
$\ell_1,\ldots,\ell_r$ of $M$; this is possible because $N_A$ is saturated,
so $N/N_A$ is free and the sequence
$0\to(N/N_A)^*\to M\to N_A^*\to0$ splits.  For $i\le k$ put
$\xi_i(x)=\langle\ell_i,x\rangle$, and write
$\xi''=(\xi_{k+1},\ldots,\xi_r)$.  The affine coordinate map
$x\mapsto(\xi_1(x),\ldots,\xi_r(x))$ has unimodular linear part.

\emph{Closedness.}  Applying $\dprime$ differentiates only the coefficient,
producing terms containing $\dprime\xi_j$ for some $j>k$; each such factor
already occurs in the wedge, so every term vanishes.  The same argument
applies to $\ddprime$.

\emph{Positivity.}  For $w>0$ the form is a nonnegative function times
$\bigwedge_j\eta_j\wedge J\eta_j$ with $\eta_j=\dprime\xi_j$, which is
strongly positive by definition.

\emph{Convergence.}  Let $\beta\in A^{k,k}_c$ and expand it in the monomials
of the basis $\xi_1,\ldots,\xi_r$.  In the product
$\alpha^A_\varepsilon\wedge\beta$, a monomial of $\beta$ containing
$\dprime\xi_j$ or $\ddprime\xi_j$ with $j>k$ is annihilated by repetition,
and a monomial omitting some $\dprime\xi_i$ or $\ddprime\xi_i$ with
$i\le k$ cannot reach bidegree $(r,r)$.  Hence only the coefficient $b$ of
$\bigwedge_{i\le k}\dprime\xi_i\wedge\ddprime\xi_i$ survives, and, the
factors $\dprime\xi\wedge\ddprime\xi$ being of even total degree and
therefore central,
\[
 \alpha^A_\varepsilon\wedge\beta
 =w\,\rho_\varepsilon(\xi'')\,b\;\tau_r .
\]
Since the linear parts $\ell_1,\ldots,\ell_r$ form a lattice basis of $M$,
the lattice measure $\omega_N$ is
$\rd\xi_1\cdots\rd\xi_r$, and Fubini gives
\[
 \int_{N_{\mathbf R}}\alpha^A_\varepsilon\wedge\beta
 =w\int_{\mathbf R^{r-k}}\rho_\varepsilon(u)
   \Bigl(\int_{\{\xi''=u\}}b\,\rd\xi_1\cdots\rd\xi_k\Bigr)\rd u .
\]
The inner integral is a continuous compactly supported function of $u$,
because $b$ is smooth with compact support; as $\rho_\varepsilon$ is an
approximate identity, the limit is its value at $u=0$, namely
$w\int_Ab\,\omega_A=\int_{(A,w)}\beta$.
\end{proof}

\subsection{Transverse intersection as a current product}
\label{sec:ch11-transverse}

\begin{definition}
\label{def:ch11-transverse-to-A}
Let $X$ be a tropical $n$-cycle in $N_{\mathbf R}$, represented by a
weighted complex $(\mathscr C,m)$, let $A$ be a rational affine subspace of
dimension $k$ with weight $w$, and let $U\subset N_{\mathbf R}$ be open.  We
say that $X$ and $(A,w)$ meet \emph{transversely over $U$} if
\begin{enumerate}
\item $|X|\cap A\cap U$ is contained in the union of the relative interiors
  of the facets of $\mathscr C$;
\item every facet $\sigma$ with $\sigma\cap A\cap U\ne\varnothing$ satisfies
  $N_{\sigma,\mathbf R}+N_{A,\mathbf R}=N_{\mathbf R}$.
\end{enumerate}
In that case the \emph{transverse intersection} $X\cdot(A,w)$ over $U$ is
the restriction to $U$ of the weighted polyhedral intersection; its local
facets are the sets $\sigma\cap A\cap U$, of pure dimension
$c=n+k-r$, with weights
\[
 m_{\sigma\cap A}=m_\sigma\,w\,[N:N_\sigma+N_A].
\]
\end{definition}

This is the transverse case of the weight formula
\eqref{eq:tropical-stable-weight}, with $Y$ the cycle $(A,w)$; condition
$(2)$ is the requirement that every facet of the intersection have the
expected dimension.  Note that $(2)$ forces $n+k\ge r$, so that
$c=n+k-r\ge0$, unless no facet meets $A$ inside $U$, in which case
$|X|\cap A\cap U=\varnothing$ and both sides of every formula below vanish.
When $c=0$ the intersection is a finite set of points with multiplicities,
and $\int_{X\cdot(A,w)}\beta$ is the corresponding weighted sum of values of
$\beta$.

\begin{theorem}
\label{thm:ch11-transverse-current-product}
Let $X=(\mathscr C,m)$ and $(A,w)$ meet transversely over an open set $U$,
and let $\alpha^A_\varepsilon$ be as in
\Cref{lem:ch11-linear-regularization}.  Then, for every
$\beta\in A^{c,c}_c(U)$ with $c=n+k-r$,
\[
 \lim_{\varepsilon\to0}\ \int_{(\mathscr C,m)}\alpha^A_\varepsilon\wedge\beta
 \;=\;\int_{X\cdot(A,w)}\beta .
\]
Equivalently,
$\delta_X\wedge[\alpha^A_\varepsilon]\to\delta_{X\cdot(A,w)}$ weakly on $U$.
\end{theorem}

\begin{proof}
\emph{Step 1: localization.}  Let $K=|X|\cap\Supp\beta$, a compact subset of
$U$.  The set $K\cap A$ is compact, and by transversality it is covered by
the relative interiors of finitely many facets $\sigma_1,\ldots,\sigma_s$
with $N_{\sigma_i,\mathbf R}+N_{A,\mathbf R}=N_{\mathbf R}$.  If some
sequence $\varepsilon_\nu\to0$ admitted points
$x_\nu\in K\cap\Supp\alpha^A_{\varepsilon_\nu}$ outside a fixed
neighbourhood $W$ of $K\cap A$, then, $K$ being compact, a subsequence would
converge to a point of $K$ at distance $0$ from $A$, hence to a point of
$K\cap A$, contradicting $x_\nu\notin W$.  Therefore, for $\varepsilon$
small, $\Supp\alpha^A_\varepsilon\cap K$ lies in $W$, and we may take $W$ so
small that it meets only the facets $\sigma_i$ and no cell of $\mathscr C$
of dimension $<n$.  Consequently
\[
 \int_{(\mathscr C,m)}\alpha^A_\varepsilon\wedge\beta
 =\sum_{i=1}^sm_{\sigma_i}\int_{\sigma_i}
   \alpha^A_\varepsilon\wedge\beta
\]
for all small $\varepsilon$.

\emph{Step 2: one facet.}  Fix $\sigma=\sigma_i$ and write
$N_\sigma\cap N_A=N_{\sigma\cap A}$, which is saturated in $N_\sigma$ and
has rank $c$, by the dimension count
$\dim(\sigma\cap A)=n+k-r$.  Choose a lattice basis $v_1,\ldots,v_n$ of
$N_\sigma$ whose first $c$ members form a lattice basis of
$N_{\sigma\cap A}$, and let $s_1,\ldots,s_n$ be the dual affine coordinates
on the span of $\sigma$, so that $\omega_\sigma=\rd s$ and
$\sigma\cap A=\{s_{c+1}=\cdots=s_n=0\}$ locally.  The restrictions of
$\xi_{k+1},\ldots,\xi_r$ to that span vanish on $N_{\sigma\cap A,\mathbf R}$
and are affine, so
\[
 \xi''|_\sigma=D\,(s_{c+1},\ldots,s_n)^{\mathsf T}+\text{const},
\]
where $D$ is the square matrix, of size $r-k=n-c$, of the induced
isomorphism
$N_\sigma/N_{\sigma\cap A}\to N/N_A$ in the chosen lattice bases.  Its
determinant satisfies
\[
 |\det D|
 =\bigl[N/N_A:\pi(N_\sigma)\bigr]
 =\bigl[N:N_\sigma+N_A\bigr],
\]
where $\pi:N\to N/N_A$ is the projection; the first equality is the
definition of the index of the image lattice, and the second follows from
$\pi(N_\sigma)=(N_\sigma+N_A)/N_A$.

Restricting the form,
\[
 \Bigl(\bigwedge_{j>k}\dprime\xi_j\wedge\ddprime\xi_j\Bigr)\Big|_\sigma
 =\det(D)^2\bigwedge_{i>c}\dprime s_i\wedge\ddprime s_i ,
\]
by the computation of \Cref{lem:ch11-orientation-drops-out} applied to $D$.
As in the proof of \Cref{lem:ch11-linear-regularization}, only the
coefficient $b$ of $\bigwedge_{i\le c}\dprime s_i\wedge\ddprime s_i$ in
$\beta|_\sigma$ survives the wedge, so
\[
 \int_\sigma\alpha^A_\varepsilon\wedge\beta
 =w\det(D)^2\int\rho_\varepsilon\bigl(\xi''(s)\bigr)\,b(s)\,\rd s .
\]
Substituting $u=\xi''(s)$ in the last $n-c$ variables, whose Jacobian is
$|\det D|$, and applying Fubini,
\[
 \begin{aligned}
  \int_\sigma\alpha^A_\varepsilon\wedge\beta
   &=w\,|\det D|\int_{\mathbf R^{n-c}}\rho_\varepsilon(u)
     \Bigl(\int b\bigl(s',s''(u)\bigr)\,\rd s'\Bigr)\rd u\\
   &\xrightarrow[\varepsilon\to0]{}
     w\,|\det D|\int_{\sigma\cap A}\beta ,
 \end{aligned}
\]
since the inner integral is continuous in $u$ with compact support and
$\int_{\sigma\cap A}\beta=\int b(s',0)\,\rd s'$ by definition of the lattice
measure of $\sigma\cap A$.

\emph{Step 3: summation.}  Multiplying by $m_\sigma$ and summing over the
facets gives
\[
 \lim_{\varepsilon\to0}\int_{(\mathscr C,m)}\alpha^A_\varepsilon\wedge\beta
 =\sum_im_{\sigma_i}w\,[N:N_{\sigma_i}+N_A]\int_{\sigma_i\cap A}\beta
 =\int_{X\cdot(A,w)}\beta ,
\]
by \Cref{def:ch11-transverse-to-A}.
\end{proof}

\begin{corollary}
\label{cor:ch11-transverse-balanced}
In the situation of \Cref{thm:ch11-transverse-current-product}, if $X$ is
balanced then so is $X\cdot(A,w)$ over $U$.
\end{corollary}

\begin{proof}
By \Cref{thm:ch11-current-closed} it suffices to show that
$\int_{X\cdot(A,w)}\dprime\eta=0$ for every
$\eta\in A^{c-1,c}_c(U)$.  By the theorem this is
$\lim_\varepsilon\int_X\alpha^A_\varepsilon\wedge\dprime\eta$.  Since
$\alpha^A_\varepsilon$ is $\dprime$-closed and of even total degree
$2(r-k)$, the Leibniz rule gives
$\dprime(\alpha^A_\varepsilon\wedge\eta)
=\alpha^A_\varepsilon\wedge\dprime\eta$, and
$\alpha^A_\varepsilon\wedge\eta$ has compact support on $|X|$ by Step 1 of
the proof.  \Cref{thm:balanced-stokes} then gives
$\int_X\alpha^A_\varepsilon\wedge\dprime\eta=0$ for every small
$\varepsilon$, whence the limit is zero.
\end{proof}

\subsection{From transverse to stable intersection}
\label{sec:ch11-stable-current}

\begin{corollary}
\label{cor:ch11-stable-current-product}
Let $X$ be a tropical $n$-cycle in $N_{\mathbf R}$ and let $(A,w)$ be a
weighted rational affine subspace of dimension $k$ with $n+k=r$, so that the
stable product $X\cdot(A,w)$ of
\Cref{thm:tropical-stable-intersection} is a $0$-cycle.  Let $v$ be generic
and let $\Omega$ be a bounded open set.  Then for all sufficiently small
$t>0$ the cycles $X$ and $(A+tv,w)$ meet transversely over $\Omega$, and
\[
 \lim_{t\downarrow0}\ \lim_{\varepsilon\to0}\
   \delta_X\wedge\bigl[\alpha^{A+tv}_\varepsilon\bigr]
 =\delta_{X\cdot(A,w)}
 \qquad\text{weakly on }\Omega .
\]
\end{corollary}

\begin{proof}
Transversality of $X$ with a generic small translate of $A$ is the
displacement statement recalled in \Cref{chap:tropical-spaces} before
\eqref{eq:tropical-stable-weight}; by
\Cref{thm:ch11-transverse-current-product} the inner limit is
$\delta_{X\cdot(A+tv,w)}$, a finite sum of weighted Dirac masses in
$\Omega$.

Let $\beta$ be a smooth function with compact support $K\Subset\Omega$.
Choose $\Omega_0$ with $K\Subset\Omega_0\Subset\Omega$.  Only finitely many
cells of $\mathscr C$ meet $\overline{\Omega_0}$, so the displacement
construction of \Cref{thm:tropical-stable-intersection} reduces there to a
finite polyhedral calculation.  Let $q_1,\ldots,q_M$ be the points of the
stable $0$-cycle $X\cdot(A,w)$ lying in $K$.  For all sufficiently small
$t>0$, the transverse intersection points
that can meet $K$ split into finite clusters
$q_{j,1}(t),\ldots,q_{j,s_j}(t)$ with $q_{j,h}(t)\to q_j$.  If
$\mu_{j,h}$ is the transverse multiplicity of $q_{j,h}(t)$, the local
displacement formula \eqref{eq:tropical-stable-weight} says
\[
 \sum_{h=1}^{s_j}\mu_{j,h}=m_{X\cdot(A,w)}(q_j).
\]
This is the point at which one must use the stable $0$-cycle rather than
the set-theoretic intersection $|X|\cap A$, which can be positive
dimensional.

Points outside these clusters remain outside $K$ for small $t$.  Therefore
\[
 \begin{aligned}
 \int_{X\cdot(A+tv,w)}\beta
 &=\sum_{j=1}^M\sum_{h=1}^{s_j}
      \mu_{j,h}\,\beta\bigl(q_{j,h}(t)\bigr)\\
 &\longrightarrow
   \sum_{j=1}^M m_{X\cdot(A,w)}(q_j)\,\beta(q_j)
  =\int_{X\cdot(A,w)}\beta,
 \end{aligned}
\]
because $\beta(q_{j,h}(t))\to\beta(q_j)$ in every finite cluster.
\end{proof}

\begin{remark}
\label{rem:ch11-stable-positive-dim}
When $c=n+k-r>0$, an analogous statement also has to control convergence of
the moving facets and their lattice measures, rather than only the locations
and weights of finitely many points.  That positive-dimensional convergence
is not proved here and is not used later in the book.
\end{remark}

\begin{example}[Two tropical lines]
\label{ex:ch11-two-lines-current}
Let $L\subset\mathbf R^2$ be the tropical line of
\Cref{fig:tropical-line-balancing} and let $A=\mathbf R(1,0)$ be the
horizontal line, of weight $1$, a rational affine subspace with
$N_A=\mathbf Z(1,0)$ and $k=1$, so that $c=1+1-2=0$.  Here $A$ contains the
horizontal ray of $L$, so the two do not meet transversely and
\Cref{cor:ch11-stable-current-product} must be used.

Displace $A$ downwards: $A_t=A-t(0,1)$ with $t>0$.  Then $A_t$ meets $|L|$
in the single point $p_t=(-t,-t)$ on the ray of direction $(-1,-1)$, and
that intersection is transverse, since
$\mathbf Z(1,1)+\mathbf Z(1,0)=\mathbf Z^2$.  Its multiplicity is
\[
 1\cdot1\cdot[\mathbf Z^2:\mathbf Z(1,1)+\mathbf Z(1,0)]
 =\bigl|\det\begin{psmallmatrix}1&1\\1&0\end{psmallmatrix}\bigr|=1 .
\]
Hence $L\cdot(A_t,1)=[p_t]$, and letting $t\downarrow0$ gives
$L\cdot(A,1)=[0]$: one point of multiplicity one at the vertex.  As
$\varepsilon\to0$ the regularized forms $\alpha^{A_t}_\varepsilon$ converge
to $\delta_{A_t}$, and the double limit is $\delta_0$, in accordance with
the stable-intersection computation of \Cref{chap:tropical-spaces}.

Displacing upwards instead, $A_t'=A+t(0,1)$ meets $|L|$ at $(0,t)$ on the
vertical ray, again with multiplicity
$|\det\begin{psmallmatrix}0&1\\1&0\end{psmallmatrix}|=1$, and the limit is
again $\delta_0$.  The independence of the displacement direction, proved in
general in \Cref{thm:tropical-stable-intersection}, is here the statement
that the two chambers give the same total weight; balancing at the vertex is
what makes them agree.
\end{example}

\begin{example}[The lattice-index test as a limit of currents]
\label{ex:ch11-index-test-current}
Let $X$ be the line $\mathbf R(1,2)\subset\mathbf R^2$ with weight $1$,
which is balanced because it has no codimension-one cell, and let
$A=\mathbf R(1,0)$ with weight $1$.  They meet transversely at the origin
and
\[
 [N:N_X+N_A]
 =\bigl|\det\begin{psmallmatrix}1&2\\1&0\end{psmallmatrix}\bigr|=2 ,
\]
so $X\cdot(A,1)=2[0]$.  Let us verify this through
\Cref{thm:ch11-transverse-current-product}.  Take $\xi_2=x_2$ as the
defining functional of $A$, so that
$\alpha^A_\varepsilon=\rho_\varepsilon(x_2)\,\dprime x_2\wedge\ddprime x_2$,
and let $\beta$ be a compactly supported function.  Parametrize $X$ by its
lattice coordinate $s\mapsto s(1,2)$; then $x_2=2s$ and
\[
 \alpha^A_\varepsilon\big|_X=\rho_\varepsilon(2s)\,
   \dprime(2s)\wedge\ddprime(2s)
 =4\rho_\varepsilon(2s)\,\dprime s\wedge\ddprime s,
\]
whence
\[
 \int_X\alpha^A_\varepsilon\wedge\beta
 =4\int_{\mathbf R}\rho_\varepsilon(2s)\beta(s(1,2))\,\rd s
 =2\int_{\mathbf R}\rho_\varepsilon(u)\beta\bigl(\tfrac u2(1,2)\bigr)\,\rd u
 \longrightarrow 2\beta(0).
\]
Once more the factor $4$ from the two labelled copies is halved by the
change of variables, and what remains is the lattice index $2$.  A
construction that produced $1$ here would be the one warned against in the
lattice-index example of \Cref{chap:tropical-spaces}.
\end{example}

\begin{remark}[What is proved here and what is not]
\label{rem:ch11-intersection-scope}
Three limitations should be stated plainly.

First, one factor is a \emph{linear} cycle, that is, a weighted rational
affine subspace.  This is the case used in
\Cref{chap:analytic-tropicalization} for dimension and multiplicity
formulas, in \Cref{chap:cld-currents} for the fibre-integration and
projection statements
\Cref{thm:cld-fibre-integration,thm:cld-fibre-projection}, and in
\Cref{chap:cld-bedford-taylor} for localization of Bedford--Taylor products
along a coordinate subspace.  The general stable product of two tropical
cycles is not treated here.

Second, even in that case we prove convergence of a regularized product,
not the existence of a product of two currents.  The functional
$\delta_X\wedge\delta_A$ is not defined in the smooth calculus, and no
choice of regularizing family should be read as defining it.  What the
theorem asserts is that the answer is independent of the mollifier $\rho$,
which is visible from the proof.

Third, the extension in which such products do exist intrinsically is the
delta-form calculus of \Cref{chap:delta-forms}, where
\Cref{thm:delta-preform-algebra} produces a bigraded algebra whose
multiplication induces stable tropical intersection on the polyhedral part.
The comparison between the limit computed here and the delta-form product
is \Cref{prop:tropical-delta-pl}; the two agree wherever both are defined,
which is why the present regularized statement suffices for the
applications.
\end{remark}

\section{Dictionary and sign ledger}
\label{sec:ch11-ledger}

\subsection{From this chapter to the general theory}
\label{sec:ch11-dict-forward}

\begin{center}
\small
\begin{tabular}{p{.30\textwidth}|p{.28\textwidth}|p{.32\textwidth}}
\textbf{Established here}&\textbf{Generalized as}&\textbf{Where}\\ \hline
lattice measure $\omega_\sigma$
  &calibration of a parlinear space
  &\cref{chap:tropical-spaces},
   \cite[\S4.6]{CLD25}\\
$\int_\sigma f\tau_n=\int_\sigma f\,\omega_\sigma$
  &integration of $(n,n)$-forms on a tropical space
  &\cref{thm:cld-global-stokes},
   \cref{chap:cld-forms}\\
boundary trace
   \eqref{eq:superform-boundary-traces}
  &boundary integral $\int^\partial$ of CLD
  &\cref{sec:appJ-currents},
   \cite[\S5.9]{CLD25}\\
boundary defect $\mathbf N(\tau)$
   (\cref{def:ch11-defect})
  &discordance of a calibration
  &\cite[\S4.6.7]{CLD25}\\
balanced $\Leftrightarrow$ Stokes
   (\cref{thm:ch11-balancing-iff-stokes})
  &closedness of an integration current
  &\cref{thm:ch11-current-closed},
   \cref{thm:current-fundamental-cycle}\\
Green formula
   (\cref{thm:ch11-green-defect})
  &Green formula on a Berkovich space
  &\cref{chap:cld-forms},
   \cref{chap:cld-bedford-taylor}\\
graph Laplacian comparison
   (\cref{thm:ch11-graph-green})
  &$\dprimeddprime$ of a model function on a curve
  &\cref{prop:skeletal-curvature-curve},
   \eqref{eq:cld-curve-graph-laplacian}\\
Dirichlet sign
   (\cref{cor:tropical-dirichlet-sign})
  &energy pairing and its negativity
  &\cref{cor:ch08-dirichlet-energy},
   \cref{chap:cld-bedford-taylor}\\
pushforward with index
   (\cref{def:ch11-proper-pushforward})
  &proper pushforward of currents
  &\cref{prop:current-pushforward-formal},
   \cref{thm:cld-fibre-projection}\\
projection formula
   (\cref{cor:ch11-projection-formula})
  &$g_*(T\wedge g^*\alpha)=g_*T\wedge\alpha$
  &\cref{prop:current-pushforward-formal}\\
transverse product as a limit
   (\cref{thm:ch11-transverse-current-product})
  &delta-form multiplication
  &\cref{thm:delta-preform-algebra},
   \cref{prop:tropical-delta-pl}
\end{tabular}
\end{center}

\subsection{The sign table}
\label{sec:ch11-sign-table}

Every convention fixed in this chapter is listed here with the test that
determines it and the place in \Cref{app:normalizations} where it is
recorded.  A reader recomputing any formula of
\Crefrange{chap:cld-forms}{chap:cld-examples} should be able to reconstruct
each row from its test.

\begin{center}
\small
\begin{tabular}{p{.26\textwidth}|p{.30\textwidth}|p{.16\textwidth}|p{.18\textwidth}}
\textbf{Convention}&\textbf{Fixed by}&\textbf{Value}&\textbf{Recorded in}\\
\hline
positive top orientation $\tau_n$
  &interlaced, not grouped
  &$\tau_n=c_n\,\dprime x_{[n]}\wedge\ddprime x_{[n]}$
  &\eqref{eq:appJ-lagerberg-orientation}\\
grouped-to-interlaced constant
  &\cref{lem:ch11-interlaced-grouped}
  &$c_p=(-1)^{p(p-1)/2}$
  &\eqref{eq:appJ-cp}\\
measure on a cell
  &covolume one for $N_\sigma$
  &$\omega_\sigma$, unsigned
  &\cref{sec:appJ-cld}\\
orientation of a cell
  &drops out of every integral
  &none needed
  &\cref{lem:ch11-orientation-drops-out}\\
contraction $\iota'',\iota'$
  &acts in one labelled block only
  &no traversal sign
  &\cref{sec:appJ-currents}\\
normal direction
  &outward, primitive in $N_\sigma/N_\tau$
  &$n_{\sigma,\tau}=-u_{\sigma/\tau}$
  &\cref{sec:appJ-currents}\\
trace prefactor, bidegree $(n-1,n)$
  &\cref{lem:ch11-trace-normal-form}
  &$(-1)^{n-1}=c_nc_{n-1}$
  &\eqref{eq:superform-boundary-traces}\\
trace prefactor, bidegree $(n,n-1)$
  &\cref{lem:ch11-trace-normal-form}
  &$-1$
  &\eqref{eq:superform-boundary-traces}\\
boundary defect
  &sum of outward normals
  &$\mathbf N(\tau)$
  &\cref{def:ch11-defect}\\
graph Laplacian versus trace
  &outgoing slopes point inward
  &$\Delta_\Gamma=-\sum\operatorname{tr}'$
  &\cref{lem:ch11-graph-trace}\\
Hessian sign
  &convex gives positive
  &$\dprimeddprime u=(u_{ij})$
  &\eqref{eq:appJ-real-hessian}\\
Dirichlet pairing
  &$\int f\dprimeddprime f\wedge\alpha\le0$
  &negative semidefinite
  &\cref{cor:tropical-dirichlet-sign}\\
pushforward multiplicity
  &one surviving lattice index
  &$[N'_{\sigma'}:F(N_\sigma)]$
  &\cref{def:ch11-proper-pushforward}\\
transverse multiplicity
  &one surviving lattice index
  &$[N:N_\sigma+N_A]$
  &\eqref{eq:tropical-stable-weight}
\end{tabular}
\end{center}

\subsection{The unit tests}
\label{sec:ch11-unit-tests}

The chapter contains the following complete numerical checks.  Together they
exercise every row of the sign table at least once.

\begin{center}
\small
\begin{tabular}{p{.34\textwidth}|p{.28\textwidth}|p{.30\textwidth}}
\textbf{Test}&\textbf{Detects}&\textbf{Where}\\ \hline
strongly positive form on
  $\mathbf R^2$ has positive integral
  &interlaced versus grouped
  &\cref{ex:ch11-positivity-sign},
   \cref{ex:ch11-grouped-wrong-sign}\\
divergence theorem on $[0,1]^2$
  &prefactor $(-1)^{n-1}$
  &\cref{ex:ch11-square-wrong-sign}\\
hypotenuse of the standard triangle
  &primitive versus Euclidean normal
  &\cref{ex:ch11-euclidean-normal-trap},
   \cref{ex:ch11-triangle}\\
$\operatorname{SL}(2,\mathbf Z)$ chart change
  &basis independence
  &\cref{ex:ch11-gl2z-test}\\
$t\mapsto2t$ on a segment
  &index appears once, not twice
  &\cref{ex:ch11-nonunimodular}\\
Stokes on a segment
  &outward normal at each end
  &\cref{ex:ch11-segment-stokes}\\
Stokes on a tripod
  &cancellation by balancing
  &\cref{ex:ch11-tripod-stokes}\\
tripod with three rays in a half-plane
  &necessity of balancing
  &\cref{ex:ch11-tripod-unbalanced}\\
two weights on crossed lines
  &weights are extra data
  &\cref{ex:ch11-weights-not-support}\\
$f(t)=t$ on $[0,\ell]$
  &Laplacian sign of
   \cref{sec:metric-graphs}
  &\cref{ex:ch11-graph-segment-check}\\
tripod potential $\varphi_D$
  &agreement with
   \cref{ex:ch08-p1-tripod}
  &\cref{ex:ch11-graph-tripod-check}\\
Green function $g_q$ on $[0,1]$
  &sign of the Dirichlet energy
  &\cref{ex:ch11-green-function-check}\\
rays $(2,1),(2,-1),(-1,0)$ to $\mathbf R$
  &projection formula and index
  &\cref{ex:ch11-projection-curve}\\
integral horizontal lines
  &necessity of properness
  &\cref{ex:ch11-pushforward-nonproper}\\
two tropical lines
  &displacement independence
  &\cref{ex:ch11-two-lines-current}\\
line of slope $2$ against a coordinate line
  &transverse lattice index
  &\cref{ex:ch11-index-test-current}
\end{tabular}
\end{center}

Four points are worth memorizing on leaving the chapter.
\begin{enumerate}
\item Orientations never survive into an integral; lattices always do.
\item The only sign created by the boundary formalism is
  $c_nc_{n-1}=(-1)^{n-1}$, and it must be tested in even dimension.
\item Outgoing slopes on a graph are inward normals; the graph Laplacian is
  minus the outward flux.
\item A lattice index enters once under pushforward and once under
  transverse intersection, never twice, because the change of variables
  removes one of the two labelled copies.
\end{enumerate}

\section*{Source notes}

The Euclidean and balanced-complex proofs above expose the mechanism hidden
inside the invariant formalism. Calibrated integration, boundary, Stokes, and
Green formulae in full parlinear generality are from
\cite[\S5]{CLD25}. Delta-current versions and tropical intersection products
are postponed to \cref{chap:delta-forms}, after ordinary currents have been defined.

\paragraph{Provenance of the individual statements.}
The half-space Stokes and Green formulae for superforms on a real vector
space, which are our \Cref{lem:ch11-segment-stokes} and
\Cref{lem:ch11-mixed-cancellation} in embryonic form, go back to Lagerberg's
original paper on superforms and appear as Lemmas~1.5.9 and~1.5.10 of
\cite{CLD25}; the mixed cancellation
$\dprime\alpha\wedge\ddprime\beta=-\ddprime\alpha\wedge\dprime\beta$ for
symmetric forms is their Lemma~1.3.8.  Our proof of it, by evaluating the
involution $J$ on a top-degree form in two ways, is the one that makes the
role of symmetry visible.  The calibrated theory on a parlinear space, the
notion of \emph{discordance} of a calibration along a codimension-one face,
and the resulting notion of a \emph{harmonious} space are
\cite[\S\S4.6.7--4.6.14]{CLD25}; our \Cref{def:ch11-defect} is the
discordance in the embedded situation, and
\Cref{thm:ch11-balancing-iff-stokes} is the observation, implicit there,
that harmoniousness is not merely sufficient but necessary.
\Cref{thm:polyhedral-stokes} and \Cref{thm:balanced-stokes} correspond to
\cite[Proposition~4.6.15]{CLD25}, and \Cref{thm:cld-global-stokes} to
\cite[Proposition~5.11.1]{CLD25}.  Their proofs there proceed by a
partition-of-unity reduction to the half-space lemma and an
$\varepsilon$-truncation near the codimension-two skeleton; we have instead
reduced to the standard simplex and transported by an affine map, which
isolates the lattice-index bookkeeping in one place and makes the
$c_nc_{n-1}$ constant explicit.  Neither route is shorter than the other; the
present one is chosen because the index identity it uses,
$d_\tau\,\langle n_{\sigma,\tau},\nu_\tau\rangle=1$ in Step~2, is the same
base-times-height identity that governs pushforward, and the reader meets it
once instead of twice.

\paragraph{Polyhedral and tropical vocabulary.}
Balancing, weights, tropical cycles, stable intersection, and the
lattice-index formula for pushforward follow
\cite[\S\S4.2--4.3 and 4.6]{MR18Tropical}, as in
\Cref{chap:tropical-spaces}; that reference also contains the combinatorial
treatment of pushforward along an arbitrary integer-affine map, which we
have replaced by the argument of
\Cref{thm:ch11-pushforward-balanced}.  The systematic development of
intersection theory on abstract tropical cycles, including pullback,
rational equivalence, and the projection formula in a form closer to
classical intersection theory, is due to Allermann and Rau; it is not needed
here and is not cited by key, since no proof in this book depends on it.

\paragraph{The comparison with metric graphs.}
\Cref{thm:ch11-graph-green} is a bookkeeping statement rather than a new
theorem, but it is the kind of statement that is usually left implicit and
is then the source of sign disputes.  The metric-graph conventions it is
matched against---outgoing slopes, the atomic Laplacian
\eqref{eq:graph-laplacian}, Green's identity \eqref{eq:graph-green}, the
negative semidefiniteness of the Dirichlet pairing
(\Cref{cor:ch08-dirichlet-energy})---are those of \Cref{sec:metric-graphs},
which follows the tropical literature represented by
\cite{MR18Tropical} and the potential theory of Berkovich curves developed
by Thuillier and, for the projective line, by Baker and Rumely.  The
identification of the curve-theoretic Laplacian with the
Chambert-Loir--Ducros operator on a skeleton is carried out in
\cite{Wanner19}, and in the recent curve-test framework in
\cite{GR25PL,GR25Classical}; a Poincar\'e lemma for the associated real
forms is \cite{Jell16}.  The reader should note that the sign in
\Cref{lem:ch11-graph-trace} is the reason those comparisons are stated with
a minus sign, and that no source is at fault: outgoing slopes and outward
normals are genuinely opposite conventions, each natural in its own context.

\paragraph{Intersection products and their limits.}
The regularization of \Cref{lem:ch11-linear-regularization} is the tropical
analogue of the classical approximation of a current of integration by
smooth bump forms; in the superform setting it appears in Lagerberg's
convolution formalism, recalled in \cite[\S1.6]{CLD25}, and in the
supercurrent framework of \cite{Ber19Super}, where the metric-dependent
supercurrent of a real submanifold is compared with calibrated tropical
currents.  Those two constructions are not the same, and
\Cref{chap:cld-currents} keeps them apart.  The full multiplicative theory,
in which stable tropical intersection becomes an honest product of
delta-forms, is the Gubler--K\"unnemann calculus of \cite{GK17,GK19}, taken
up in \Cref{chap:delta-forms}; \Cref{rem:ch11-intersection-scope} records
exactly which of its statements the present chapter does and does not
anticipate.  The enlargement of the chart category by harmonic
tropicalizations is \cite{GJR26}, and the comparison of real and
non-Archimedean Monge--Amp\`ere operators in the toric case, which uses the
integration theory of this chapter, is \cite{Vil20}; the global invariant
tropical version is \cite{BGJKK25}.

\paragraph{What is new here, and what is not.}
No theorem of this chapter is new.  What is not standard is the
organization: the boundary defect is introduced before the balanced case
rather than after it, so that
\Cref{thm:ch11-defect-stokes} covers balanced and unbalanced complexes
uniformly and the balancing hypothesis appears exactly where it is used;
the converse implication of \Cref{thm:ch11-balancing-iff-stokes} is proved
rather than asserted; balancing of a pushforward is deduced from that
converse instead of from a combinatorial index computation; and every
convention is tested numerically in \Cref{sec:ch11-unit-tests}.  Readers
comparing with \cite[\S\S4.6 and 5.9--5.11]{CLD25} should expect the same
theorems in a different order and with different proofs of the elementary
steps.
